\documentclass[12pt]{amsart}
\usepackage{ amsmath, amsthm, amsfonts, amssymb, color}
 \usepackage{mathrsfs}
\usepackage{amsfonts, amsmath}
 \usepackage{amsmath,amstext,amsthm,amssymb,amsxtra}
 \usepackage{txfonts} %also pxfonts
 \usepackage[colorlinks, citecolor=blue,pagebackref,hypertexnames=false]{hyperref}
 \allowdisplaybreaks
 \usepackage{pgf,tikz}
 \usepackage{multirow}%newtablepackage1
 \usepackage{diagbox} %newtablepackage2
\usepackage{cite}
 \usepackage{tikz}

\usetikzlibrary{decorations.pathreplacing}

 %\usepackage{showkeys}
 \textwidth =166mm
 \textheight =232mm
\marginparsep=0cm
\oddsidemargin=2mm
\evensidemargin=2mm
\headheight=13pt
\headsep=0.8cm
\parskip=0pt
\hfuzz=6pt
\widowpenalty=10000
 \setlength{\topmargin}{-0.2cm}

\DeclareMathOperator*{\essinf}{ess\ inf}
\DeclareMathOperator*{\esssup}{ess\ sup}
\setcounter {tocdepth}{1}
\begin{document}

 \baselineskip 16.6pt
\hfuzz=6pt

\widowpenalty=10000

\newtheorem{cl}{Claim}
\newtheorem{theorem}{Theorem}[section]
\newtheorem{proposition}[theorem]{Proposition}
\newtheorem{coro}[theorem]{Corollary}
\newtheorem{lemma}[theorem]{Lemma}
\newtheorem{definition}[theorem]{Definition}
\newtheorem{nota}[theorem]{Notation}
%\theoremstyle{definition}
\newtheorem{assum}{Assumption}[section]
\newtheorem{example}[theorem]{Example}
\newtheorem{remark}[theorem]{Remark}
\renewcommand{\theequation}
{\thesection.\arabic{equation}}

\def\SL{\sqrt H}

\newcommand{\mar}[1]{{\marginpar{\sffamily{\scriptsize
        #1}}}}

\newcommand{\as}[1]{{\mar{AS:#1}}}
\newcommand{\cE}{\mathcal{E}}
\newcommand\R{\mathbb{R}}
\newcommand\RR{\mathbb{R}}
\newcommand\CC{\mathbb{C}}
\newcommand\D {{\rm Dunkl}}
\newcommand\NN{\mathbb{N}}
\newcommand\N {{\bf N}}
\newcommand\ZZ{\mathbb{Z}}
\newcommand\HH{\mathbb{H}}
\newcommand\Z{\mathbb{Z}}
\def\RN {\mathbb{R}^n}
\renewcommand\Re{\operatorname{Re}}
\renewcommand\Im{\operatorname{Im}}

\newcommand{\mc}{\mathcal}

\def\hs{\hspace{0.33cm}}
\newcommand{\la}{\alpha}
\newcommand{\eps}{\tau}
\newcommand{\pl}{\partial}
\newcommand{\supp}{{\rm supp}{\hspace{.05cm}}}
\newcommand{\x}{\times}
\newcommand{\lag}{\langle}
\newcommand{\rag}{\rangle}

\newcommand\wrt{\,{\rm d}}

\newcommand{\norm}[2]{|#1|_{#2}}
\newcommand{\Norm}[2]{\|#1\|_{#2}}

\title[Schatten class properties of Dunkl--Riesz transform commutators]{Weyl--Chamber geometry, Poincar\'e inequality meets Schatten class properties of Dunkl--Riesz transform commutators}



\author{Zhijie Fan}
\address{Zhijie Fan, School of Mathematical Sciences, South China Normal University, Guangzhou 510631, China}
\email{fanzhj3@mail2.sysu.edu.cn}







\author{Ji Li}
\address{Ji Li, Department of Mathematics, Macquarie University,NSW2109, Australia}
\email{ji.li@mq.edu.au}



\author{Dmitriy Zanin}
\address{Dmitriy Zanin, University of New South Wales, Kensington, NSW 2052, Australia}
\email{d.zanin@unsw.edu.au}






  \date{\today}

 \subjclass[2010]{47B10, 42B20, 43A85}
\keywords{Schatten class, Riesz transform commutators, Dunkl operator, Besov space, Nearly weakly orthogonal sequence}




\begin{abstract}

\end{abstract}

\maketitle


\tableofcontents

%\newpage
\section{Introduction}
\begin{definition}{\rm\label{beso}
Suppose $1< p< \infty$. We say that a  function $f\in L^p_{{\rm loc}}(\mathbb{R}^{N},dw)$ belongs to {Dunkl} Besov space $B^{p}_{\D}(\mathbb{R}^N)$ if
\begin{align*}
\|f\|_{B^{p}_{\D}(\mathbb{R}^N)}:=\left(\int_{\mathbb{R}^{N}}\int_{\mathbb{R}^{N}}\left(\frac{d(x,y)}{\|x-y\|}\right)^{p}\frac{|f(x)-f(y)|^p}{w(B(x,d(x,y)))^2}dw(x)dw(y)\right)^{1/p}<+\infty.
\end{align*}}
\end{definition}

\subsection{Notation}
Conventionally, we let $\mathbb{N}$ denote the set of positive integers and set $\mathbb{Z}_+ := \{0\} \cup \mathbb{N}$. Throughout the paper, the symbols $C$ and $c$ denote positive constants (possibly varying from line to line) that may depend only on $N$, $p$, and the structural data $(R,k)$ of the Dunkl setting, but are independent of the main parameters (e.g., the symbol $b$ and auxiliary decay factors).
 We write $A \lesssim B$ to mean that $A \leq C B$ for some constant $C > 0$, and $A \sim B$ to mean that both $A \lesssim B$ and $B \lesssim A$ hold.
For \(\lambda>0\) and a cube \(Q\) with sidelength \(\ell(Q)\), write \(\lambda Q\) for the cube concentric with \(Q\) and sidelength \(\lambda\,\ell(Q)\). Similarly, for every ball \(B=B(x,r)\), set \(\lambda B:=B(x,\lambda r)\). For any $1\leq p\leq\infty$, we denote by $p'$ the conjugate of $p$, which satisfies $\frac{1}{p}+\frac{1}{p'}=1$. We denote the average of a function $f$ over a $\mu$-measurable set $E$ by
\begin{align}\label{defaverage}
(f)_E:=\fint_{E}f(x)d\mu(x):=\frac{1}{\mu(E)}\int_{E}f(x)d\mu(x).
\end{align}
For a subset $E \subset \mathbb{R}^{N}$, we denote by $\chi_E$ its characteristic function.

Denote by $\{e_j\}_{j=1}^N$ the standard orthonormal basis on $\mathbb{R}^N$.
%For two measurable Borel sets $E$ and $F$, we use the notation ${\rm dist}(E,\,F)$ to denote the distance between them.


\section{Preliminaries}
\setcounter{equation}{0}
\subsection{Harmonic analysis in the Dunkl setting}
Let $\mathbb{R}^N$ be the Euclidean space equipped with the usual scalar product $\langle \cdot,\,\cdot\rangle $ and the induced norm $\|\cdot\|$. For a nonzero vector $v\in \mathbb{R}^N$, the reflection $\sigma_v$ with respect to the hyperplane $v^{\perp}$ orthogonal to $v$ is defined by
\begin{align*}
	\sigma_v x:= x-2\frac{\langle x,v\rangle}{\|v\|^2}v.
\end{align*}
A finite set $R\subset \mathbb{R}^N\backslash \{0\}$ is called a root system if $\sigma_v(R)=R$ for every $v\in R$. A root system $R$ is called a normalized reduced root system if  $\|v\|^2=2$ for every $v\in R$.
The finite group $G$ generated by the reflections $\{\sigma_v\}_{v\in R}$ is called the Weyl group (reflection group) of the root system $R$.
Let $k: R \to [0,\infty)$ be a multiplicity function, which is a $G$-invariant function  fixed throughout the paper. We will denote by $\mathcal O(x)=\{\sigma(x): \sigma \in G\}$ the $G$-orbit of a point $x\in\mathbb{R}^N$, and $\mathcal O(E)=\bigcup_{x \in E} \mathcal O(x)$ the $G$-orbit of $E\subset \mathbb{R}^N$.

Define the Dunkl measure $dw$ on $\mathbb{R}^N$ by
$$
dw(x) :=\prod_{v \in R} |\langle x,v\rangle|^{k(v)} dx,
$$
where $dx$ stands for the Lebesgue measure on $\mathbb{R}^N$. Let $B(x,r):=\{y\in \mathbb{R}^N:\|x-y\|< r\}$ be the Euclidean ball centred at $x$ with radius $r>0$. Let $\N:=N+\sum_{v\in R}k(v)$. Note that
$$w(B(tx,tr))=t^{\N}w(B(x,r)),\ {\rm for}\ x\in\mathbb{R}^N,\ t,r>0.$$
Furthermore,
\begin{align}\label{twosideinequ}
  w(B(x,r))\sim r^N \prod_{v \in R}(|\langle x, v \rangle |+r)^{k(v)},
\end{align}
where the implicit constant only depends on $N,R$ and $k$.  It follows from \eqref{twosideinequ} that there exists a constant $C\geq 1$ such that for every $x\in \mathbb{R}^N$,
\begin{align}\label{doublinginequality}
C^{-1}\left(\frac{r_2}{r_1}\right)^N\leq \frac{w(B(x,r_2))}{w(B(x,r_1))}\leq C\left(\frac{r_2}{r_1}\right)^\N,\ {\rm for}\ 0<r_1<r_2,
\end{align}
so the numbers $\N$ and $N$ are called the upper and lower homogeneous dimensions, respectively.

The Dunkl operators $T_\xi$ are defined by
$$T_\xi f(x):=\partial_\xi f(x)+\sum_{v\in R}\frac{k(v)}{2}\langle v,\xi\rangle\frac{f(x)-f(\sigma_v(x))}{\langle v,x\rangle},$$
where $\partial_\xi f$ denotes the directional derivative of $f$ in the direction $\xi$. Furthermore, the Dunkl Laplacian is defined by $$\Delta_{\D}f(x):=\sum_{j=1}^N T_{e_j}^2f(x)=\Delta_{\mathbb{R}^N}f(x)+\sum_{v\in R}k(v)\delta_v f(x),$$
where $\Delta_{\mathbb{R}^N}$ is the Laplacian on $\mathbb{R}^N$ and $\delta_v$ is defined by
$$\delta_v f(x):=\frac{\partial_v f(x)}{\langle v,x\rangle}-\frac{f(x)-f(\sigma_v(x))}{\langle v,x\rangle^2}.$$
In this article, we consider the Dunkl-Riesz transforms defined by $R_{\D,j}:=T_{e_j}(-\Delta_{\D})^{-1/2}$, $j=1,2,\ldots,N$.


\begin{definition}{\rm
The Weyl chambers are defined as the closures of connected components of the set
  \begin{align}
    \left\{x\in\mathbb{R}^N:\langle x,v\rangle \neq 0\ \ { for} \ { all} \ v\in R\right\}.\end{align}}
\end{definition}
Define $
	d(x,y):=\min_{\sigma\in G}\|x-\sigma(y)\|
$
as the distance between two $G$-orbits $\mathcal O(x)$ and $\mathcal O(y)$. Recall from \cite[Chapter VII, proof of Theorem 2.12]{MR514561} that
\begin{align}\label{distanceequa}
d(x,y)=\|x-y\|\ {\rm if}\ {\rm and}\ {\rm only}\ {\rm if}\ x,y\in\Gamma
\end{align}
for some closed Weyl chamber $\Gamma$.

It is clear that $\mathcal O(B(x,r)) = \{y\in \mathbb{R}^N: d(x,y)<r\}$. Moreover,
\[
w(B(x,r)) \leq w(\mathcal O(B(x,r))) \leq |G|w(B(x,r))
\]
for all $x \in \mathbb{R}^N$ and $r>0$. We also denote $B^d(x,r):= \{y\in \mathbb{R}^N: d(x,y)<r\}$, i.e., $B^d(x,r)=\mathcal O(B(x,r))$.

Denote by $K_{\D,j}(x,y)$ the integral kernel of $R_{\D,j}$. We recall the kernel estimates for the Riesz transforms in the Dunkl setting proved in \cite{MR4733961}. Although these estimates may not be sharp (see the comment following Theorem~2.1 in \cite[p.~10]{MR4733961}), they suffice for our purposes.
\begin{lemma}\label{kernelestimate}
There exists a constant $C$ such that for $j\in\{1,2,\ldots,N\}$ and every $x,y$ with $d(x,y)\neq 0$,
\begin{align*}
|K_{\D,j}(x,y)|\leq C\frac{d(x,y)}{\|x-y\|}\frac{1}{w(B(x,d(x,y)))}.
\end{align*}
\end{lemma}
\begin{proof}
The proof is given in \cite[Theorem 1.1]{MR4733961}.
\end{proof}

\begin{lemma}\label{nondegenerate}
For $j\in\{1,2,\ldots,N\}$, and for every ball $B=B(x_0,r)\subset \mathbb{R}^N$, there is another ball $\widehat{B}=B(y_0,r)$ with $y_0=x_0+5re_j$ such that for every $(x,y)\in B\times \widehat{B}$, $K_{\D,j}(x,y)$ does not change sign and satisfies
\begin{align*}
|K_{\D,j}(x,y)|\geq\frac{C}{w(B(x_0,r))}.
\end{align*}
\end{lemma}
\begin{proof}
The lower bound is given in \cite[Theorem 1.2]{MR4733961}, while the fact that $K_{\D,j}(x,y)$ does not change sign can be seen in the proof of \cite[Theorem 1.2]{MR4733961}. Indeed, with the choice $y_0:=x_0+5re_j$, we have $y_j-x_j\geq 3r$ for all $(x,y)\in B\times \widehat{B}$. Together with the fact that the heat kernel $h_t(x,y)$ associated with the Dunkl Laplacian is  positive, the representation
$$K_{\D,j}(x,y)=-\frac{C}{2}\int_0^\infty \frac{y_j-x_j}{t}h_t(x,y)\frac{dt}{\sqrt{t}},$$
implies that $K_{\D,j}(x,y)<0$ (and thus, does not change sign) on $ B\times \widehat{B}$.
\end{proof}



\subsection{Systems of dyadic cubes}


\begin{definition}\label{Defsystem}{\rm
Let $X\subseteq \mathbb{R}^N$ be a nonempty open set. A countable family
$
    \mathscr{D}
    := \bigcup_{k\in\mathbb{Z}}\mathscr{D}_{k}, \
    \mathscr{D}_{k}
    :=\{Q^k_{\alpha}\colon \alpha\in \mathscr{A}_k\},
$
of Borel sets $Q^k_{\alpha}\subseteq X$ is called \textit{a
system of dyadic cubes over  $X$}  if it has the following properties:

\begin{enumerate}
  \item[(i)] For each $k\in\mathbb{Z}$, we have $$    X
    = \bigcup_{\alpha\in \mathscr{A}_k} Q^k_{\alpha};$$\smallskip

  \item[(ii)] $\text{If }\ell\geq k\text{, then either }
        Q^{\ell}_{\beta}\subseteq Q^k_{\alpha}\text{ or }
        Q^k_{\alpha}\cap Q^{\ell}_{\beta}=\varnothing$;\smallskip

  \item[(iii)] $
    \text{For each }(k,\alpha)\text{ and each } \ell\leq k,
    \text{ there exists a unique } \beta
    \text{ such that }Q^{k}_{\alpha}\subseteq Q^\ell_{\beta};
$

\smallskip

\item[(iv)] For each $(k,\alpha)$ there exist $2^N$ indices $\beta$ such that
    $$ Q^{k+1}_{\beta}\subseteq Q^k_{\alpha},\ {\rm and}\
        Q^k_{\alpha} =\bigcup_{{\substack{Q\in\mathscr{D}_{k+1}\\ Q\subseteq Q^k_{\alpha}}}}Q;$$

\smallskip

\item[(v)] For each $(k,\alpha)$, one has $$B_X(x^k_{\alpha},A_1 2^{-k})
    \subseteq Q^k_{\alpha}\subseteq B_X(x^k_{\alpha},A_2 2^{-k})
    =: B(Q^k_{\alpha})$$
    for some positive constants $A_1,A_2$ depending only on the geometry of $X$, {where we set $B_X(x,r):=B(x,r)\cap X$ for all $x\in X$ and $r>0$.}

\smallskip

\item [(vi)]  If $\ell\geq k$ and
   $Q^{\ell}_{\beta}\subseteq Q^k_{\alpha}$, then
   $$B(Q^{\ell}_{\beta})\subseteq B(Q^k_{\alpha}).$$
\end{enumerate}



   }
\end{definition}

Each $Q^k_{\alpha}$ is called a \textit{dyadic cube of generation} $k$, with center $x^k_{\alpha}\in Q^k_{\alpha}$ and side-length $\ell(Q^k_{\alpha})=2^{-k}$.
\begin{nota}
For every dyadic cube $Q\in\mathscr{D}$, we denote by $c_Q$ and $\ell(Q)$ its center and side-length, respectively.
\end{nota}

\begin{nota}\label{standardsystem}
Let $$\mathbb{D}^0:=\bigcup_{k\in\mathbb{Z}}\mathbb{D}_k^0$$ be the standard system of dyadic cubes over $\mathbb{R}^N$.
Here each dyadic cube is a product of half-open intervals
\[
[m_1 2^{-k}, (m_1+1)2^{-k}) \times \cdots \times [m_N 2^{-k}, (m_N+1)2^{-k}),
\]
with $(m_1,\ldots,m_N)\in \mathbb{Z}^N$. Moreover, we let $$\mathbb{D}_{\Gamma}:=\bigcup_{k\in\mathbb{Z}}\mathbb{D}_{\Gamma,k}$$ be the system of dyadic cubes over a fixed chamber $\Gamma$, whose  construction will be given in Corollary \ref{Construction over gamma}.
\end{nota}


\begin{lemma}\label{Nsystems}
There exist $2^N$ systems of dyadic cubes $\mathbb{D}^1,\mathbb{D}^2,\ldots,\mathbb{D}^{2^N}$ such that for any ball $B=B(x_B,r_B)\subset \mathbb{R}^N$ with $\frac{2^{-k-2}}{3}\leq r_B< \frac{2^{-k-1}}{3}$ for some $k\in\mathbb{Z}$,  there exists a cube ${\displaystyle Q\in\bigcup_{\nu=1}^{2^N}\mathbb{D}^\nu_k}$ such that
$$B\subset Q\subset 12\sqrt{N} B.$$
\end{lemma}

\begin{proof}
{\bf $\bullet$ Step 1.}  {\it Recall the one-dimensional setting given in \cite[Section 2]{MR3420475}.}

Let $\mathbb{D}^{(0)}(\mathbb{R})$ be the standard system of dyadic intervals over $\mathbb{R}$ (see Notation \ref{standardsystem}).
Let $\delta=1/3$ and define
$$
   \mathbb D^{(\delta)}(\mathbb R)=\bigcup_{k\in\mathbb Z} \mathbb D^{(\delta)}_k(\mathbb R),
$$
where for $k\geq0$,
$$
    \mathbb D^{(\delta)}_k(\mathbb R)=\{I+\delta \mid I\in \mathbb D_k^{(0)}(\mathbb R)\},
$$
while for $k<0$ and $k$ even,
$$
    \mathbb D^{(\delta)}_k(\mathbb R)
    =\Big\{\Big[{m\over 2^k}+\delta+\sum_{j=(k/2)+1}^0 2^{-2j},\ {m+1\over 2^k}
    +\delta+\sum_{j=(k/2)+1}^0 2^{-2j}\Big)\ \Big|\ m\in
    \mathbb Z\Big\}.
$$
These choices together with the nestedness property property completely
determine the collections $\mathbb D^{(\delta)}_k(\mathbb R)$ for $k<0$, $k$ odd.

Fix an interval $J\subset\mathbb R$ with $2^{-k-1}/3\le |J|<2^{-k}/3$ for some $k\in\mathbb Z$. Now we set
$A_k=\{ m\cdot 2^{-k} \mid m\in\mathbb Z \}$ for every fixed $k$ and
\begin{enumerate}
    \item[(1)]$A_k^{(\delta)}=\{ \delta+ m\cdot 2^{-k} \mid m\in\mathbb Z
        \}$ for $k\geq 0$,

    \item[(2)]$A_k^{(\delta)}=\{ \delta+
        \sum_{j=(k/2)+1}^{0}2^{-2j}+ m\cdot 2^{-k} \mid m\in\mathbb Z
        \}$ for $k<0$, $k$ even, and

    \item[(3)]$A_k^{(\delta)}=\{ \delta+
        \sum_{j=(k-1)/2+1}^{0}2^{-2j}+ m\cdot 2^{-k} \mid
        m\in\mathbb Z \}$ for $k<0$, $k$ odd.
\end{enumerate}
Note that for any two different points $a,b\in A_k\cup
A_k^{(\delta)}$, we have $|a-b|\geq 2^{-k}/3>|J|$. Thus,
there is at most one element of $A_k\cup A_k^{(\delta)}$ contained in
$J$. Hence, $J\cap A_k=\emptyset $ or $J\cap
A_k^{(\delta)}=\emptyset $. Therefore, $J$ must be contained in
some dyadic interval $I\in \mathbb D_k^{(0)}$ or $I\in \mathbb D^{(\delta)}_k$ and
$|I|=2^{-k}\leq 6 |J|$.

{\bf $\bullet$ Step 2.} {\it We extend the construction to $\mathbb{R}^N$ ($N\geq 1$).}

For $\vartheta=(\vartheta_1,\vartheta_2,\ldots,\vartheta_N)\in\{0,\delta\}^N$, define
$$\mathbb{D}^{(\vartheta)}:=\bigcup_{k\in\mathbb{Z}}\mathbb{D}_k^{(\vartheta)},\quad \mathbb{D}_k^{(\vartheta)}:=\left\{\prod_{j=1}^N I_j: I_j\in \mathbb{D}_k^{(\vartheta_j)}\right\}.$$
Applying Step 1 to $\pi_i(B)$, where $\pi_i(B)$ denotes the $i$-th coordinate projection, we obtain $I_i\in\mathbb{D}_k^{(0)}$ or $I_i\in\mathbb{D}_k^{(\delta)}$ with $\pi_i(B)\subset I_i$ and $|I_i|=2^{-k}\leq 12 r_B$. Let $\vartheta_i\in\{0,\delta\}$
be the corresponding shift so that $I_i\in\mathbb D_k^{(\vartheta_i)}$, set $\vartheta:=(\vartheta_1,\ldots,\vartheta_N)$, and let $Q=\prod_{j=1}^NI_j\in\mathbb{D}_k^{(\vartheta)}$. It follows that $B\subset Q$ and $\ell(Q)=2^{-k}\leq 12r_B$. Hence, for every $y\in Q$,
$$\|y-x_B\|\leq \sqrt{N}\ell(Q)\leq 12\sqrt{N}r_B.$$
This implies that $Q\subset B(x_B,12\sqrt{N}r_B)$. Enumerating $\big\{\mathbb D^{(\vartheta)}:\vartheta\in\{0,\delta\}^N\big\}$ as
$\mathbb D^1,\ldots,\mathbb D^{2^N}$, we complete the proof of Lemma \ref{Nsystems}.
\end{proof}



%Hence, for any ball $B\subset \mathbb{R}^N$, via reducing to each axis in $\mathbb{R}^N$
%and repeating the argument in $\mathbb{R}$, we see that there exist $k\in\mathbb Z$ and a cube ${\displaystyle Q\in\bigcup_{\nu=1}^{2^N}\mathbb{D}^\nu_k}$ such that
%$$B\subset Q\subset \kappa_N B,$$
%where $\kappa_N$ is a positive constant depending only on $N$.


\subsection{Haar Basis}\label{HaarSection}



Given a system of dyadic cubes $\mathscr{D}$ over $\mathbb{R}^N$, we now recall the standard way to construct an orthonormal Haar system corresponding to $\mathscr{D}$. For each $k\in\mathbb{Z}$ and $Q\in\mathscr{D}_k$, denote by $Ch(Q):=\{R\in\mathscr{D}_{k+1}:R\subseteq Q\}$. Moreover,  we let
$$h_Q^0:=w(Q)^{-1/2}\chi_Q.$$
Fix an enumeration $\{Q_i\}_{i=1}^{2^N}$ of $Ch(Q)$. Let
$$\hat{Q}_i:=\bigcup_{j=i}^{2^N}Q_j$$
and
$$h_Q^i:=\sqrt{\frac{w(Q_i)w(\hat{Q}_{i+1})}{w(\hat{Q}_i)}}\Bigg(\frac{\chi_{Q_i}}{w(Q_i)}-\frac{\chi_{\hat{Q}_{i+1}}}{w(\hat{Q}_{i+1})}\Bigg),\quad i=1,2,\ldots,2^N-1.$$


\begin{lemma}\label{Haarexpansion}
For any $1<p<\infty$ and $f\in
    L^p(\mathbb{R}^N,dw)$, we have
    \[
        f(x)
        =  \sum_{Q\in\mathscr{D}}\sum_{i=1}^{2^N-1}
            \langle f,h^{i}_{Q}\rangle_{L^2(\mathbb{R}^N,dw)} h^{i}_{Q}(x), %\quad \text{if } \mu(X)=\infty,
    \]
    where the sum converges both in the
    $L^p(\mathbb{R}^N,dw)$ norm and pointwise almost everywhere.
 \end{lemma}


\begin{lemma}
The family of Haar functions $\{h_{Q}^{i}\}_{Q\in\mathscr{D},0\leq i\leq 2^N-1}$ has the following properties:
    \begin{itemize}
\item[(i)] $h_{Q}^{i}$ is supported on $Q$ for $i\in\{0,1,\ldots,2^N-1\}$;
\item[(ii)] $h_{Q}^{i}$ is constant on each $R\in Ch(Q)$ for $i\in\{0,1,\ldots,2^N-1\}$;
\item[(iii)] $\int_{\mathbb{R}^N} h_{Q}^{i}(x)dw(x)= 0$, for $i\in\{1,2,\ldots,2^N-1\}$;
\item[(iv)] $\langle h_{Q}^{i},h_{Q}^{j}\rangle_{L^2(\mathbb{R}^N,dw)} = 0$ for
$i \neq j$, $i$, $j\in\{1, 2,\ldots, 2^N- 1\}$;
\item[(v)] The collection
$\{h_{Q}^{i}\}_{i = 0}^{2^N-1}$
is an orthogonal basis for the vector space $V(Q)$ of all functions on $Q$ that are constant on each sub-cube $R\in Ch(Q)$.
\item[(vi)]
            $
 \Norm{h_{Q}^{i}}{L^p(\mathbb{R}^N,dw)}
    \approx w(Q)^{\frac{1}{p} - \frac{1}{2}}$
   \quad for~1 $\leq p \leq \infty$ and $i\in\{1,2,\ldots,2^N-1\}$;

    \end{itemize}
\end{lemma}


\subsection{Nearly weakly  orthonormal sequences}
Rochberg and Semmes \cite{RS} introduced the notion of \emph{nearly weakly  orthonormal (NWO)} sequences of functions. This notion has since proved highly effective in characterizing Schatten class membership of  singular integral commutators in a variety of contexts (see \cite{FLLXarxiv,HKarxiv,MR4756021,MR4552581,MR4873796,LLW,RS,Zhangwei1,Zhangwei2}). The properties of NWO sequences were originally investigated in the Euclidean setting (equipped with Lebesgue measure), but the extension to spaces of homogeneous type is straightforward once a system of dyadic cubes is available in this general framework (see \cite[Section 13]{HKarxiv} for a complete proof).  For our purposes, we adapt  the extension to the Dunkl setting $(\mathbb{R}^N,\|\cdot\|,w)$ rather than recalling the general  definition
of NWO sequences and the general extension to spaces of homogeneous type.

{\it Let $\mathscr{D}$ be a system of dyadic cubes. A sequence of functions $\{e_Q\}_{Q\in\mathscr{D}}$ is an NWO sequence if
\begin{enumerate}
  \item $\supp e_Q\subseteq cQ$ for some positive constant $c$ independent of $Q$
  \item $\|e_Q\|_{L^s(\mathbb{R}^{N},dw)}\leq w(Q)^{1/s-1/2}$ for some $s>2$.
\end{enumerate}
}
\begin{lemma}\label{NWOLEMMAA}
Let $1<p<\infty$, and let $\mathscr{D}$ be a
system of dyadic cubes over $\mathbb{R}^N$. For any compact operator $ T $ on $ L ^2 (\mathbb R ^{N},dw)$, we have
\begin{equation*}
\Bigl(
\sum_{Q\in \mathscr{D}} |\langle T e_Q, f_Q \rangle_{L^2(\mathbb{R}^N,dw)} | ^{p}
\Bigr) ^{1/p} \leq C_{N,p,w} \| T \| _{S ^{p}},
\end{equation*}
where $\{e_Q\}_{Q\in\mathscr{D}}$ and $\{f_Q\}_{Q\in\mathscr{D}}$ are NWO sequences.
\end{lemma}
\begin{proof}
This property can be found in  \cite[(1.10)]{RS} (see also \cite[Section 7]{HKarxiv}).
\end{proof}

\begin{lemma}\label{NWOLEMMAA2}
Let $0<p<\infty$, $0<q\leq \infty$, and let $\mathscr{D}$ be a
system of dyadic cubes over $\mathbb{R}^N$. Assume that
$$A=\sum_{Q\in\mathscr{D}}\lambda_Q\langle \cdot,\,e_Q\rangle_{L^2(\mathbb{R}^N,dw)} f_Q,$$
where $\{e_Q\}_{Q\in\mathscr{D}}$ and $\{f_Q\}_{Q\in\mathscr{D}}$ are NWO sequences and $\{\lambda_Q\}_{Q\in\mathscr{D}}$ is a sequence of scalars. We have
\begin{align}\label{nwose}
\|A\|_{S^{p,q}}\leq C_{N,p,q,w} \|\{\lambda_Q\}_Q\|_{\ell^{p,q}}.
\end{align}
\end{lemma}
\begin{proof}
This property can be found in  \cite[(1.9)]{RS} (see also \cite[Section 7]{HKarxiv}).
\end{proof}

\section{Analytic tools for Weyl chamber geometry}
\setcounter{equation}{0}

{\bf Convention. }
Throughout this paper, we will use the following notation:
\begin{itemize}
  \item Let $R
  \subset \mathbb{R}^{N}$ be a finite normalized reduced root system;
  \item Let
$W:=\operatorname{span}(R)$ with $N_{0}:=\dim W$;
\item Let
$
  \Omega:=\mathbb{R}_+^{N_{0}}\times \mathbb{R}^{N-N_{0}};
$
\item Let $\Gamma_+$ be a fixed open Weyl chamber (see Definition \ref{def:weyl_chamber_positive_roots}) and $\Gamma$ be its closure.
\end{itemize}


\begin{definition}{\rm \cite[Definition ~2.7]{HeckmanCoxeterGroups}}\label{def:weyl_chamber_positive_roots} A connected component of \begin{align}
    \left\{x\in\mathbb{R}^N:\langle x,v\rangle \neq 0\ \ { for} \ { all} \ v\in R\right\}
  \end{align} is called an open
Weyl chamber. Its closure is called a Weyl chamber.  Define the corresponding sets of \emph{positive} and \emph{negative} roots by
\[
R_{+}:=\{\alpha\in R:\ \langle x,\alpha\rangle>0\ \text{for all }x\in \Gamma_{+}\},
\qquad
R_{-}:=-R_{+}.
\]
\end{definition}

\begin{definition}{\rm \cite[Definition~2.9]{HeckmanCoxeterGroups}}\label{def:simple_roots}
A root $\alpha\in R_{+}$
is called simple in $R_{+}$ if it cannot be written as
\[
\alpha=\lambda_{1}\alpha_{1}+\lambda_{2}\alpha_{2}
\quad\text{with }\lambda_{1},\lambda_{2}\ge 1\ \text{and}\ \alpha_{1},\alpha_{2}\in R_{+}.
\]
The set $\Delta\subset R_{+}$ of all simple roots is called the
simple system.
\end{definition}

\begin{lemma}\label{lem:simple_system_associated_to_chamber}
With the notation given in Definition \ref{def:simple_roots},  $\Delta$ is a basis of $W$. If we denote by
$
\Delta=\{\alpha_{1},\alpha_{2},\ldots,\alpha_{N_{0}}\},
$
then the open and closed chambers can be written as
\begin{align}
\Gamma_{+}
&=\Bigl\{x\in\mathbb{R}^{N}:\ \langle x,\alpha_{i}\rangle>0,\ i=1,\ldots,N_{0}\Bigr\},
\label{eq:open_chamber_by_simple_roots}\\
\Gamma
&=\Bigl\{x\in\mathbb{R}^{N}:\ \langle x,\alpha_{i}\rangle\ge 0,\ i=1,\ldots,N_{0}\Bigr\},
\label{eq:closed_chamber_by_simple_roots}
\end{align}
and every root $\alpha\in R$ admits a unique expansion
\[
\alpha=\sum_{i=1}^{N_{0}}c_{i}\alpha_{i}
\]
such that either $c_{i}\ge 0$ for all $i$ (when $\alpha\in R_{+}$), or $c_{i}\le 0$ for all $i$
(when $\alpha\in R_{-}$). In particular, the simple system $\Delta$ is uniquely determined by $\Gamma_{+}$
(equivalently, by $\Gamma$).
\end{lemma}

\begin{proof}
%We regard $R$ as a normalized root system in the Euclidean space $W$ endowed with the
%restricted inner product. Since $\sigma_{\alpha}(\alpha)=-\alpha$ and $\sigma_{\alpha}(R)=R$
%for every $\alpha\in R$, we have $R=-R$.

Let $\Delta=\{\alpha_{1},\ldots,\alpha_{n}\}\subset R_{+}$ be the set of simple roots in $R_{+}$. By \cite[Proposition~2.10]{HeckmanCoxeterGroups}, every $\alpha\in R_{+}$
is a nonnegative linear combination of $\alpha_{1},\ldots,\alpha_{n}$. Since $R=R_{+}\cup R_{-}$ (this follows from the fact that for each $\alpha\in R$, the continuous function
$x\mapsto\langle x,\alpha\rangle$ does not vanish on the connected set $\Gamma_{+}$) and $R_{-}=-R_{+}$, it follows that
$R\subset\operatorname{span}(\Delta)$ and hence,
\[
W=\operatorname{span}(R)\subset\operatorname{span}(\Delta).
\]
The reverse inclusion $\operatorname{span}(\Delta)\subset W$ is immediate from $\Delta\subset R\subset W$,
so $\operatorname{span}(\Delta)=W$. On the other hand, it follows from \cite[Theorem~2.16]{HeckmanCoxeterGroups} that the simple roots are linearly
independent. Consequently, $\Delta$ is a basis
of $W$, and therefore $n=\dim W=N_{0}$. Moreover, every $\alpha\in R$ has a unique expansion
$\alpha=\sum_{i=1}^{N_{0}}c_{i}\alpha_{i}$. If $\alpha\in R_{+}$, the coefficients can be chosen
nonnegative by \cite[Proposition~2.10]{HeckmanCoxeterGroups}. If $\alpha\in R_{-}=-R_{+}$, then
$-\alpha\in R_{+}$ has an expansion with all coefficients nonnegative, and multiplying by $-1$ gives
an expansion of $\alpha$ with all coefficients nonpositive.
The descriptions \eqref{eq:open_chamber_by_simple_roots} and \eqref{eq:closed_chamber_by_simple_roots} follow from \cite[Corollary~2.11]{HeckmanCoxeterGroups}.


To see that $\Delta$ is uniquely determined by $\Gamma_+$, note that $\Gamma_+$ determines $R_{+}$
by Definition~\ref{def:weyl_chamber_positive_roots}, and the set of simple roots in $R_{+}$ is
characterized intrinsically by Definition~\ref{def:simple_roots}. Finally, since $\Gamma_+$ is the (unique) interior component of $\Gamma$, specifying $\Gamma$ determines $\Gamma_+$ and vice versa. This completes the proof of Lemma \ref{lem:simple_system_associated_to_chamber}.
\end{proof}
\begin{nota}
With the notation given in Lemma \ref{lem:simple_system_associated_to_chamber}, denote by
 $$\Delta=\{\alpha_{1},\ldots,\alpha_{N_{0}}\}$$ the simple system associated with $\Gamma$.
\end{nota}

\begin{lemma}\label{chambergeometry}
Each Weyl chamber $\Gamma$ admits the decomposition
\[
\Gamma = K(\Gamma)\times W^\perp,
\qquad
K(\Gamma):=\{w\in W:\ \langle w,\alpha\rangle\ge 0\ \text{for all }\alpha\in\Delta\}.
\]
\end{lemma}

\begin{proof}
Write $x=w+z$ with $w\in W$ and $z\in W^{\perp}$. Since $\alpha_{i}\in W$, we have $\langle x,\alpha_{i}\rangle=\langle w,\alpha_{i}\rangle$ for all
$i=1,\ldots,N_0$. This, together with \eqref{eq:closed_chamber_by_simple_roots}, implies that $x\in \Gamma$ if and only if
$
w\in K(\Gamma).
$
This implies that
$
\Gamma
=
K(\Gamma)\times W^{\perp},
$
and thus completes the proof.
\end{proof}

%
%\begin{coro}
%Let $R\subset\mathbb{R}^N$ be a finite normalized reduced root system, let $W:=\operatorname{span}(R)$ and $N_{0}:=\dim W$, and let $D$ be any Weyl chamber in $\mathbb{R}^N$. Then there exists a linear isomorphism
%\[
%\Phi:\mathbb{R}^N\longrightarrow \mathbb{R}^N
%\]
%such that
%\[
%\Phi(D)=\mathbb{R}_+^{\,N_{0}}\times \mathbb{R}^{\,N-N_{0}}.
%\]
%\end{coro}
%
%\begin{proof}
%By the theorem \ref{chambergeometry}, we have the product decomposition
%\[
%D=K(D)\times W^\perp,
%\qquad
%K(D)=\{w\in W:\ \langle w,\alpha\rangle\ge0\ \ \forall\,\alpha\in\Delta\},
%\]
%where $\Delta=\Delta(D)=\{\alpha_1,\dots,\alpha_{N_{0}}\}\subset R$ consists of $N_{0}$ linearly independent facet normals; in particular, $K(D)$ is a simplicial cone in $W$ given by the intersection of exactly $N_{0}:=\dim W$ closed half-spaces with linearly independent outward normals $\Delta$.
%
%Define $L:W\to\mathbb{R}^{N_{0}}$ by
%\[
%L(w):=\big(\langle w,\alpha_1\rangle,\dots,\langle w,\alpha_{N_{0}}\rangle\big).
%\]
%Relative to the basis $\Delta$ of $W$, the matrix of $L$ is the Gram matrix
%$G=(\langle \alpha_j,\alpha_i\rangle)_{i,j}$, which is nonsingular because
%$\Delta$ is linearly independent. Hence $L$ is a linear isomorphism $W\!\xrightarrow{\;\sim\;}\mathbb{R}^{N_{0}}$. From the description of $K(D)$ it follows that
%\[
%L\big(K(D)\big)=\{t\in\mathbb{R}^{N_{0}}:\ t_i\ge0\ \text{for all }i\}=\mathbb{R}_+^{\,N_{0}}.
%\]
%Choose any linear isomorphism $U:W^\perp\!\xrightarrow{\;\sim\;}\mathbb{R}^{N-N_{0}}$, and define
%\[
%\Phi:\mathbb{R}^N=W\oplus W^\perp\longrightarrow \mathbb{R}^{N_{0}}\oplus \mathbb{R}^{N-N_{0}}\cong\mathbb{R}^N,
%\qquad
%\Phi(w+z):=\big(L(w),\,U(z)\big).
%\]
%Then $\Phi$ is a linear isomorphism and
%\[
%\Phi(D)=\Phi\big(K(D)\times W^\perp\big)=L\big(K(D)\big)\times U\big(W^\perp\big)
%=\mathbb{R}_+^{\,N_{0}}\times \mathbb{R}^{\,N-N_{0}},
%\]
%as claimed.
%\end{proof}



\begin{coro}\label{coro:chamber-orthant}
Let $\Phi_{W^{\perp}}$ be any linear isomorphism from $W^\perp$ to $\mathbb{R}^{N-N_{0}}$, and let
\[
\Phi_W:W\longrightarrow\mathbb{R}^{N_{0}},
\qquad
\Phi_W(w):=\big(\langle w,\alpha_1\rangle,\dots,\langle w,\alpha_{N_{0}}\rangle\big),
\]
\[
\Phi:\mathbb{R}^N=W\oplus W^\perp\longrightarrow
\mathbb{R}^{N_{0}}\oplus\mathbb{R}^{N-N_{0}}\cong\mathbb{R}^N,
\qquad
\Phi(w+z):=\big(\Phi_W(w),\,\Phi_{W^\perp}(z)\big).
\]
Then $\Phi$ is a linear isomorphism and
\[
\Phi(\Gamma)=\mathbb{R}_+^{N_{0}}\times \mathbb{R}^{N-N_{0}}.
\]
\end{coro}

\begin{proof}
By Lemma~\ref{chambergeometry} we have $\Gamma=K(\Gamma)\times W^\perp$, where
\[
K(\Gamma)=\{w\in W:\ \langle w,\alpha_i\rangle\ge0,\ i=1,\dots,N_{0}\}.
\]
Since $\Delta$ is a basis of $W$, the Gram matrix
$G=(\langle \alpha_i,\alpha_j\rangle)_{i,j=1}^{N_{0}}$ is nonsingular. Hence,
$\Phi_W:W\longrightarrow\mathbb{R}^{N_{0}}$ is a linear isomorphism. Moreover,
for $w\in W$ we have $w\in K(\Gamma)$ if and only if $\Phi_{W}(w)\in\mathbb{R}_+^{N_{0}}$,
so
\[
\Phi_{W}\big(K(\Gamma)\big)=\mathbb{R}_+^{N_{0}}.
\]
By construction, $\Phi$ is a linear isomorphism (as a direct sum of isomorphisms),
and therefore
\[
\Phi(\Gamma)
=\Phi\big(K(\Gamma)\times W^\perp\big)
=\Phi_{W}\big(K(\Gamma)\big)\times \Phi_{W^{\perp}}\big(W^\perp\big)
=\mathbb{R}_+^{N_{0}}\times \mathbb{R}^{N-N_{0}},
\]
which proves the desired conclusion.
\end{proof}


\begin{coro}\label{Construction over gamma}
 Let $\Gamma$ be an arbitrary Weyl chamber. There exists a system of dyadic cubes over $\Gamma$.
\end{coro}\begin{proof}
Let $\Phi:\mathbb{R}^N\to\mathbb{R}^N$ be the linear isomorphism from Corollary~\ref{coro:chamber-orthant} so that
\[
\Phi(\Gamma)=\Omega:=\mathbb{R}_+^{N_0}\times \mathbb{R}^{N-N_0}.
\]

For each $k\in\mathbb{Z}$, define $\mathbb{D}^0_{\Omega,k}$ to be the family of half-open cubes
\[
Q=\prod_{i=1}^{N} I_i,
\qquad
I_i=\big[m_i2^{-k},(m_i+1)2^{-k}\big),
\]
where $m_i\in\mathbb{Z}_+$ for $1\le i\le N_0$ and $m_i\in\mathbb{Z}$ for $N_0<i\le N$. Set
\[
\mathbb{D}^0_\Omega:=\bigcup_{k\in\mathbb{Z}}\mathbb{D}^0_{\Omega,k}.
\]
It is standard that $\mathbb{D}^0_\Omega$ is a system of dyadic cubes over $\Omega$.

Next, for each $k\in\mathbb{Z}$, define
\[
\mathbb{D}_{\Gamma,k}:=\big\{\Phi^{-1}(Q): Q\in \mathbb{D}^0_{\Omega,k}\big\},
\qquad
\mathbb{D}_{\Gamma}:=\bigcup_{k\in\mathbb{Z}}\mathbb{D}_{\Gamma,k}.
\]
Since $\Phi:\mathbb{R}^N\to\mathbb{R}^N$ is a linear isomorphism, we see that $\mathbb{D}_{\Gamma}$ is a system of dyadic cubes over $\Gamma$.
\end{proof}




\begin{nota}
 Let $\Phi$ be the linear isomorphism given in Corollary~\ref{coro:chamber-orthant}. For every Borel set $E\subset\Omega$, define
\begin{align}\label{defmu}
  \mu(E):=w\big(\Phi^{-1}(E)\big).
\end{align}
\end{nota}
\begin{lemma}\label{mulemma}
We have
  \begin{equation}\label{pushforward-density}
d\mu(u,z)=\phi(u)\,du\,dz,\quad \phi(u)=|{\rm det}(\Phi)|^{-1}\prod_{v\in R_+}\Big(\sum_{i=1}^{N_{0}} a_{v,i}\,u_i\Big)^{2k(v)}.
\end{equation}
where for each $v\in R_+$, the coefficients $a_{v,i}\ge0$ are uniquely determined by the simple-root expansion
\(v=\sum_{i=1}^{N_0} a_{v,i}\alpha_i\) and satisfy $(a_{v,1},\dots,a_{v,N_0})\neq\mathbf 0$.
\end{lemma}
\begin{proof}
Since $\Phi(\Gamma)=\Omega$, we have $\Phi^{-1}(\Omega)=\Gamma$. Hence, for every nonnegative Borel function
$f$ on $\Omega$,
\[
\int_{\Omega} f(u,z)\,d\mu(u,z)
=\int_{\Gamma} f\big(\Phi(x)\big)\,dw(x).
\]
We make the linear change of variables $(u,z)=\Phi(x)$, so that $x=\Phi^{-1}(u,z)$. Since $\Phi$ is linear and invertible, the Jacobian is constant and
\[
dx = |{\rm det}(\Phi)|^{-1}\,du\,dz.
\]
It follows from the definition of $\Phi$ that
$
u_i=\langle x,\alpha_i\rangle$ for $i=1,\dots,N_0.
$
This implies that
\begin{align}\label{changevariable}
\langle x,v\rangle
=\Big\langle x,\sum_{i=1}^{N_0} a_{v,i}\alpha_i\Big\rangle
=\sum_{i=1}^{N_0} a_{v,i}\langle x,\alpha_i\rangle
=\sum_{i=1}^{N_0} a_{v,i}u_i.
\end{align}
Substituting these identities into the previous integral, we obtain
\begin{align*}
\int_{\Omega} f(u,z)\,d\mu(u,z)
&=\int_{\Gamma} f\big(\Phi(x)\big)\Big(\prod_{v\in R_+}|\langle x,v\rangle|^{2k(v)}\Big)\,dx \\
&=\int_{\Omega} f(u,z)\,
|{\rm det}(\Phi)|^{-1}\prod_{v\in R_+}\Big(\sum_{i=1}^{N_0} a_{v,i}u_i\Big)^{2k(v)}\,du\,dz.
\end{align*}
Therefore, \eqref{pushforward-density} holds. This completes the proof of Lemma~\ref{mulemma}.
\end{proof}


\begin{lemma}\label{lem:gallery-simple}
Let $\Gamma_1$ and $\Gamma_2$ be two Weyl chambers. Then there exist an integer $M\ge 0$ and
$\alpha_1,\dots,\alpha_M\in R$ such that, if we set $\Gamma^{(0)}:=\Gamma_2$ and
\[
\Gamma^{(k)}:=\sigma_{\alpha_k}\big(\Gamma^{(k-1)}\big),\qquad k=1,\dots,M,
\]
then the following hold:
\begin{enumerate}
  \item\label{geo1} $\alpha_k\in \Delta(\Gamma^{(k-1)})$ for each $k=1,\dots,M$;
  \item\label{geo2} $\Gamma^{(M)}=\Gamma_1$.
\end{enumerate}
\end{lemma}
\begin{proof}
If $\Gamma_1=\Gamma_2$, we take $M=0$ and there is nothing to prove. Assume $\Gamma_1\neq \Gamma_2$.
Since the Weyl group $G$ acts transitively on the set of Weyl chambers (see \cite[Theorem 2.12 (2)]{HeckmanCoxeterGroups}), there exists
$g\in G$ such that $\Gamma_1=g(\Gamma_2)$.

Let $\Delta(\Gamma_2)=\{\beta_1,\beta_2,\ldots,\beta_{N_0}\}$ be the simple system associated with $\Gamma_2$.
By \cite[Theorem 2.12(4)]{HeckmanCoxeterGroups}, the Weyl group $G$ is generated by the simple reflections $\sigma_{\beta_1},\sigma_{\beta_2},\ldots,\sigma_{\beta_{N_0}}$. Hence, there exist indices
$i_1,\ldots,i_M\in\{1,\ldots,N_0\}$ such that
\[
g=\sigma_{\beta_{i_1}}\sigma_{\beta_{i_2}}\cdots \sigma_{\beta_{i_M}}.
\]
Define
\[
w_0:=\mathrm{Id},\qquad w_k:=\sigma_{\beta_{i_1}}\cdots \sigma_{\beta_{i_k}}\quad (k=1,\ldots,M),
\]
and set
\[
\Gamma^{(k)}:=w_k(\Gamma_2)\qquad (k=0,1,\ldots,M).
\]
It follows that $\Gamma^{(0)}=\Gamma_2$ and $\Gamma^{(M)}=w_M(\Gamma_2)=g(\Gamma_2)=\Gamma_1$, which implies the statement \eqref{geo2}.

For each $k=1,\ldots,M$, define
\[
\alpha_k:=w_{k-1}(\beta_{i_k})\in R.
\]
The rest of the proof is divided into two steps.

{\bf $\bullet$ Step 1.} {\it Verify that $\Gamma^{(k)}=\sigma_{\alpha_k}(\Gamma^{(k-1)})$.}

Since each $w\in G$ is orthogonal, for any root $\beta\in R$ one has the conjugation identity
\[
w\,\sigma_{\beta}\,w^{-1}=\sigma_{w(\beta)},\qquad w\in G.
\]
Applying this with $w=w_{k-1}$ and $\beta=\beta_{i_k}$ yields
\[
\sigma_{\alpha_k}=\sigma_{w_{k-1}(\beta_{i_k})}=w_{k-1}\sigma_{\beta_{i_k}}w_{k-1}^{-1}.
\]
Therefore,
\[
\sigma_{\alpha_k}(\Gamma^{(k-1)})
= w_{k-1}\sigma_{\beta_{i_k}}w_{k-1}^{-1}\big(w_{k-1}(\Gamma_2)\big)
= w_{k-1}\sigma_{\beta_{i_k}}(\Gamma_2)
= w_k(\Gamma_2)
= \Gamma^{(k)}.
\]

{\bf $\bullet$ Step 2.} {\it Verify that $\alpha_k\in\Delta(\Gamma^{(k-1)})$.}

Recall that if $\Gamma$ is a Weyl chamber associated with simple system $\Delta(\Gamma)$, then
\[
\Gamma=\{x\in\mathbb R^N:\ \langle x,\gamma\rangle\ge 0\ \text{for all }\gamma\in\Delta(\Gamma)\}.
\]
Applying $w_{k-1}\in G$ to the above description for $\Gamma_2$ and using
$\langle w_{k-1}x,\,\beta\rangle=\langle x,\,w_{k-1}^{-1}\beta\rangle$, we obtain
\[
\Gamma^{(k-1)}=w_{k-1}(\Gamma_2)
=\Bigl\{x\in\mathbb R^N:\ \langle x,\,w_{k-1}(\beta_j)\rangle\ge 0\ \text{for }j=1,\ldots,N_0\Bigr\}.
\]
By the uniqueness of the simple system associated with a Weyl chamber
(see Lemma~\ref{lem:simple_system_associated_to_chamber}), we conclude that
\[
\Delta(\Gamma^{(k-1)})=\Delta\big(w_{k-1}(\Gamma_2)\big)=w_{k-1}\big(\Delta(\Gamma_2)\big)=\{w_{k-1}(\beta_1),\ldots,w_{k-1}(\beta_{N_0})\}.
\]
In particular, $\alpha_k=w_{k-1}(\beta_{i_k})\in\Delta(\Gamma^{(k-1)})$, proving the first statement \eqref{geo1}.

This completes the proof of Lemma \ref{lem:gallery-simple}.
\end{proof}


\begin{definition}
We say that $R'\subset\Omega$ is an axis--parallel rectangle if
\[
R' := I_1\times\cdots\times I_N,
\]
with $I_i=[I_i^-,I_i^+]\subset[0,\infty)$ for $1\le i\le N_{0}$ and $I_i=[I_i^-,I_i^+]\subset\mathbb{R}$ for $i>N_{0}$.
\end{definition}

Denote by
$$B_\Omega(x_0,r):=\{x\in\Omega:\ \|x-x_0\|<r\}.$$
The following lemma demonstrates that $(\Omega,\|\cdot\|,d\mu)$ is a doubling metric measure space.
\begin{lemma}\label{lem:ball-2sided-N}
There exists a constant $C\ge 1$ such that
for every $x_0=(u_0,z_0)\in\Omega$ and every $r>0$,
\begin{equation}\label{eq:ball-2sided-Omega}
C^{-1}\,r^N\prod_{v\in R_+}\Big(L_{v,0}+r\Big)^{2k(v)}
\ \le\
\mu\big(B_\Omega(x_0,r)\big)
\ \le\
C\,r^N\prod_{v\in R_+}\Big(L_{v,0}+r\Big)^{2k(v)},
\end{equation}
where
\[
L_v(u):=\sum_{i=1}^{N_{0}} a_{v,i}u_i,
\qquad
L_{v,0}:=L_v(u_0)=\sum_{i=1}^{N_{0}} a_{v,i}(u_0)_i.
\]
In particular, there exists a constant $C_{\mathrm{dbl},\Omega}\geq 1$ such that for every $x\in\Omega$ and every $r>0$,
\[
\mu\big(B_\Omega(x,2r)\big)\ \le\ C_{\mathrm{dbl},\Omega}\;\mu\big(B_\Omega(x,r)\big).
\]
\end{lemma}

\begin{proof}
Fix $x_0=(u_0,z_0)\in\Omega$ and $r>0$. For each $v\in R_+$ define the linear form
\[
L_v(u):=a_v\cdot u,\qquad a_v:=(a_{v,1},\dots,a_{v,N_{0}})\in \mathbb{R}_+^{N_{0}}\setminus\{0\}.
\]
Since $R_+$ is finite and each $a_v\neq 0$, we have
\[
L_*:=\max_{v\in R_+}\|a_v\|<\infty,
\qquad
m_*:=\min_{v\in R_+}\sum_{i=1}^{N_{0}} a_{v,i}>0.
\]

\noindent
\textbf{Upper bound.}
Let $x=(u,z)\in B_\Omega(x_0,r)$. Then
\[
\|u-u_0\|\ \le\ \|x-x_0\|\ \le\ r.
\]
Hence, for every $v\in R_+$,
\[
|L_v(u)-L_{v,0}|
=\big|a_v\cdot (u-u_0)\big|
\le \|a_v\|\cdot\|u-u_0\|
\le L_*\,r.
\]
Since $L_v(u)\ge 0$ on $\Omega$ (all coefficients and coordinates are nonnegative), we obtain
\[
L_v(u)\le L_{v,0}+L_*r\le C_1\,(L_{v,0}+r),
\qquad
C_1:=\max\{1,L_*\}.
\]
Therefore,
\[
\phi(u)
=\prod_{v\in R_+} L_v(u)^{2k(v)}
\le C_1^{2\sum_{v\in R_+}k(v)}\prod_{v\in R_+}(L_{v,0}+r)^{2k(v)}
=:C_2\prod_{v\in R_+}(L_{v,0}+r)^{2k(v)}.
\]
Integrating over $B_\Omega(x_0,r)$ and using
$|B_\Omega(x_0,r)|\le |B(x_0,r)|=\omega_N r^N$, where $\omega_N$ is the volume of the unit ball in $\mathbb{R}^N$, gives
\[
\mu\big(B_\Omega(x_0,r)\big)
=\int_{B_\Omega(x_0,r)}\phi(u)\,dx
\le C_2\,|B(x_0,r)|\,\prod_{v\in R_+}(L_{v,0}+r)^{2k(v)}
\le C_3\,r^N\prod_{v\in R_+}(L_{v,0}+r)^{2k(v)},
\]
with $C_3:=C_2\,\omega_N$.

\medskip\noindent
\textbf{Lower bound.}
We construct a smaller Euclidean ball $B_c$ contained in $B_\Omega(x_0,r)$ on which $\phi$ admits a uniform
positive lower bound in terms of $\prod_{v\in R_+}(L_{v,0}+r)^{2k(v)}$. To this end, we first choose
\[
\lambda:=1+\frac{L_*}{m_*}>1,
\qquad
\eta:=\frac{1}{4\big(\sqrt{N_{0}}\,\lambda+1\big)},
\qquad
\theta:=\lambda\,\eta.
\]
These parameters depend only on $L_*$ and $m_*$, and satisfy
\[
\theta> \eta>0,
\qquad
\sqrt{N_{0}}\,\theta+\eta=\eta(\sqrt{N_{0}}\lambda+1)=\frac14<1.
\]
Let $\mathbf{1}_{N_{0}}:=(1,\dots,1)\in\mathbb{R}^{N_{0}}$ and define the center
\[
x_c:=(u_c,z_c):=\big(u_0+\theta r\,\mathbf{1}_{N_{0}},\ z_0\big),
\qquad
B_c:=B(x_c,\eta r).
\]
We {\it claim} that $B_c\subset B_\Omega(x_0,r)$. To see this inclusion, we first apply the triangle inequality to see that for any $x\in B_c$,
\[
\|x-x_0\|\le \|x-x_c\|+\|x_c-x_0\|
\le \eta r+\theta r\sqrt{N_{0}}
\le r,
\]
since $\sqrt{N_{0}}\theta+\eta<1$. Hence $B_c\subset B(x_0,r)$. Moreover, for any $x=(u,z)\in B_c$, we have $\|u-u_c\|\le \|x-x_c\|\le \eta r$. This implies that for each $1\le i\le N_{0}$,
\[
u_i\ge (u_c)_i-\eta r=(u_0)_i+(\theta-\eta)r\ge 0,
\]
since $u_0\in\mathbb{R}_+^{N_{0}}$ and $\theta\ge\eta$. Thus $u\in\mathbb{R}_+^{N_{0}}$ and $x\in\Omega$. Therefore,  $B_c\subset B_\Omega(x_0,r)$. This finishes  the proof of the  claim.

\smallskip
Next we obtain a uniform lower bound for $L_v(u)$ on $B_c$.
Write any $x=(u,z)\in B_c$ as
\[
u=u_0+\theta r\,\mathbf{1}_{N_{0}}+\Delta_u,
\qquad
z=z_0+\Delta_z,
\qquad
\|\Delta_u\|^2+\|\Delta_z\|^2\le (\eta r)^2,
\]
so in particular $\|\Delta_u\|\le \eta r$.
Then for all $x=(u,z)\in B_c$ and all $v\in R_+$,
\begin{align*}
L_v(u)
&=a_v\cdot u
=a_v\cdot u_0+\theta r\,a_v\cdot\mathbf{1}_{N_{0}}+a_v\cdot\Delta_u \\
&\ge L_{v,0}+\theta r\sum_{i=1}^{N_{0}} a_{v,i}-\|a_v\|\,\|\Delta_u\|
\ge L_{v,0}+rK_0\geq c_0(L_{v,0}+r).
\end{align*}
with
\[
K_0:=\theta m_*-\eta L_*=\eta(\lambda m_*-L_*)=\eta m_*>0,\qquad c_0:=\min\{1,K_0\}\in(0,1].
\]
Consequently, for all $u\in B_c$,
\[
\phi(u)
=\prod_{v\in R_+}L_v(u)^{2k(v)}
\ge c_0^{2\sum_{v\in R_+}k(v)}\prod_{v\in R_+}(L_{v,0}+r)^{2k(v)}
=:C_4\prod_{v\in R_+}(L_{v,0}+r)^{2k(v)}.
\]
This, in combination with the fact that $B_c\subset B_\Omega(x_0,r)$, yields
\[
\mu\big(B_\Omega(x_0,r)\big)
\ge \mu(B_c)
=\int_{B_c}\phi(u)\,dx
\ge C_4\,|B_c|\prod_{v\in R_+}(L_{v,0}+r)^{2k(v)}
= C_5\,r^N\prod_{v\in R_+}(L_{v,0}+r)^{2k(v)},
\]
where $C_5:=C_4\,\omega_N\eta^N>0$.

The second statement follows immediately from the first one. Indeed, by the first statement,
\[
\mu\big(B_\Omega(x_0,2r)\big)
\le C_3\,(2r)^N\prod_{v\in R_+}(L_{v,0}+2r)^{2k(v)}
\le C_3\,2^{N+\mathcal{K}}\,r^N\prod_{v\in R_+}(L_{v,0}+r)^{2k(v)}
\le C_3C_5^{-1}\,2^{N+\mathcal{K}}\,\mu\big(B_\Omega(x_0,r)\big),
\]
where we set $\mathcal{K}:=2\sum_{v\in R_+}k(v)$. Therefore, $(\Omega,\|\cdot\|,\mu)$ is doubling with the doubling constant $C_{\mathrm{dbl},\Omega}:=C_3C_5^{-1}\,2^{N+\mathcal{K}}$.
 This finishes the proof of Lemma \ref{lem:ball-2sided-N}.
\end{proof}
Let $\Phi:\mathbb{R}^N\longrightarrow\mathbb{R}^N$ be the linear isomorphism given in  Corollary~\ref{coro:chamber-orthant} such that
$
\Phi(\Gamma)=\Omega.
$
Since $\Phi$ is linear and invertible, there exist positive constants $c_\Phi$ and $C_\Phi$ such that for any $x,y\in\mathbb{R}^N$,
\begin{equation}\label{eq:biLip-N}
c_\Phi\,\|x-y\|\ \le\ \|\Phi(x)-\Phi(y)\|\ \le\ C_\Phi\,\|x-y\|.
\end{equation}
Denote by $B_\Gamma(x,r):=\{y\in\Gamma:\ \|y-x\|<r\}$. It follows from \eqref{eq:biLip-N} that for any $x_0\in\Gamma$ and $r>0$,
\begin{align}\label{beforedoublingball}
B_{\Omega}(\Phi(x_0),c_\Phi r)\subset \Phi(B_{\Gamma}(x_0,r))\subset B_{\Omega}(\Phi(x_0),C_\Phi r).
\end{align}
This, together with Lemma  \ref{lem:ball-2sided-N}, yields
\begin{align}\label{doublingball}
  \mu(\Phi(B_{\Gamma}(x_0,r))) \simeq \mu(B_{\Omega}(\Phi(x_0),C_\Phi r)).
\end{align}
\begin{lemma}\label{construction-rectangle}
  For all $x_0\in\Gamma$ and $r>0$, one can find an axis--parallel rectangle $R'$ on $\Omega$ such that
\begin{align}\label{step3goal}
B_{\Omega}(\Phi(x_0),C_\Phi r)\subset R'\subset B_{\Omega}(\Phi(x_0),\sqrt{N}C_\Phi r).
\end{align}
Furthermore, we have
\begin{align}\label{Rdoubling}
  \mu(R')\simeq \mu(B_{\Omega}(\Phi(x_0),C_\Phi r)).
\end{align}
%\begin{align}\label{step3goal}
%\Phi(B_\Gamma(x_0,r))\subset R'\quad \text{and}\quad
%\mu(\Phi(B_\Gamma(x_0,r)))\simeq \mu(R').
%\end{align}
\end{lemma}
\begin{proof}
We write
$
(u_0,z_0):=\Phi(x_0)\in\Omega.
$
and define the axis--parallel rectangle
\begin{equation}\label{eq:rectangle-choice}
R':=\Big(\prod_{i=1}^{N_{0}} \big[\max\{0,(u_0)_i-C_\Phi r\},\ (u_0)_i+C_\Phi r\big]\Big)
\times
\Big(\prod_{j=1}^{N-N_{0}}\big[(z_0)_j-C_\Phi r,\ (z_0)_j+C_\Phi r\big]\Big).
\end{equation}

On the one hand, it follows immediately that $B_\Omega(\Phi(x_0),C_\Phi r)\subset R'$. On the other hand, since $R'$ is a rectangle with half-side lengths bounded by $C_\Phi r$ in each coordinate, we have
\begin{align*}
R'\subset B_\Omega(\Phi(x_0),\sqrt{N}\,C_\Phi r).
\end{align*}
The second statement follows from the first one together with Lemma \ref{lem:ball-2sided-N}. The proof of Lemma \ref{construction-rectangle} is finished.
\end{proof}
The following theorem demonstrates that $(\Gamma,\|\cdot\|,dw)$ is a doubling metric measure space.


\begin{theorem}\label{main:dunkldoubling}
There exists a constant $C\ge 1$ such that
for every $x\in\Gamma$ and every $r>0$,
\begin{equation}\label{eq:ball-2sided-Gamma}
C^{-1}\,r^N \prod_{v \in R_+}(|\langle x, v \rangle |+r)^{2k(v)}
\ \le\
w\big(B_\Gamma(x,r)\big)
\ \le\
C\,r^N \prod_{v \in R_+}(|\langle x, v \rangle |+r)^{2k(v)}.
\end{equation}
In particular, there exists a constant $C_{\mathrm{dbl},\Gamma}\geq 1$ such that for every $x\in\Gamma$ and every $r>0$,
\[
w\big(B_\Gamma(x,2r)\big)\ \le\ C_{\mathrm{dbl},\Gamma}\;w\big(B_\Gamma(x,r)\big).
\]
Moreover, for all $x\in\Gamma$ and $r>0$, we have
\begin{align}\label{globallocalvol}
w(B_\Gamma(x,r))\sim w(B(x,r)).
\end{align}
\end{theorem}
\begin{proof}
The main strategy of the proof is to transfer the ball estimates from $(\Omega,\|\cdot\|,\mu)$ back to $(\Gamma,\|\cdot\|,dw)$. Fix $x\in\Gamma$ and $r>0$. It follows from the equality \eqref{defmu} that
\begin{equation}\label{eq:pushforward-ball}
w\big(B_\Gamma(x,r)\big)=\mu\big(\Phi(B_\Gamma(x,r))\big).
\end{equation}
Applying $\mu$ to the inclusions \eqref{beforedoublingball} and using \eqref{eq:pushforward-ball} yields
\begin{equation}\label{eq:wball-between-muballs}
\mu\big(B_{\Omega}(\Phi(x),c_\Phi r)\big)
\ \le\
w\big(B_\Gamma(x,r)\big)
\ \le\
\mu\big(B_{\Omega}(\Phi(x),C_\Phi r)\big).
\end{equation}
Write $\Phi(x)=(u,z)\in\Omega$. For each $v\in R_+$ recall the linear form
$L_{v,0}:=\sum_{i=1}^{N_0}a_{v,i}u_i$.
Combining the inequality \eqref{eq:wball-between-muballs} with Lemma~\ref{lem:ball-2sided-N} and the equality \eqref{changevariable}, we obtain that
\begin{equation}\label{eq:wball-in-terms-of-L}
w\big(B_\Gamma(x,r)\big)\ \simeq\ r^N\prod_{v\in R_+}\big(L_{v,0}+r\big)^{2k(v)}=r^N\prod_{v\in R_+}\big(|\langle x,v\rangle|+r\big)^{2k(v)}.
\end{equation}
This proves \eqref{eq:ball-2sided-Gamma}. The second statement follows immediately from \eqref{eq:ball-2sided-Gamma}. The third statement follows from the inequalities \eqref{eq:ball-2sided-Gamma} and \eqref{twosideinequ}. This completes the proof of Theorem~\ref{main:dunkldoubling}.
\end{proof}

%\section{Test Section}
%\begin{lemma}\label{lem:Cwikel-Dunkl}
%Assume that
%\[
%N>2
%\qquad\text{and}\qquad
%\operatorname{span}(R)=\mathbb R^N.
%\]
%Then there exists a constant $C_{N,R,k}>0$ such that, for every
%$f\in L^N(\mathbb R^N,dx)$,
%\[
%M_f(-\Delta_{\D})^{-1/2}
%\in
%S^{N,\infty}\bigl(L^2(\mathbb R^N,dw)\bigr)
%\]
%and
%\begin{equation}\label{eq:Cwikel-Dunkl}
%\bigl\|
%M_f(-\Delta_{\D})^{-1/2}
%\bigr\|_
%{S^{N,\infty}(L^2(\mathbb R^N,dw))}
%\leq
%C_{N,R,k}
%\|f\|_{L^N(\mathbb R^N,dx)}.
%\end{equation}
%\end{lemma}
%
%\begin{proof}
%Fix once and for all
%\[
%r:=\frac{N+2}{2}.
%\]
%Since $N>2$, we have
%\begin{equation}\label{eq:r-range-Cwikel}
%2<r<N.
%\end{equation}
%Let
%\[
%r'=\frac{r}{r-1}.
%\]
%Then $r'<2$.
%
%We divide the proof into five steps.
%
%\medskip
%\noindent
%\textbf{Step 1. A localized version of the rank-one decomposition.}
%
%Let $K_{\D}(x,y)$ denote the integral kernel of
%$(-\Delta_{\D})^{-1/2}$.
%We first record a refinement of the kernel estimate used in the proof of
%Lemma~\ref{lem:rank-one-decomposition-Dunkl}.
%
%Recall that
%\[
%V(x,y,t)
%:=
%\max\bigl\{
%w(B(x,t)),w(B(y,t))
%\bigr\}.
%\]
%The mixed derivative estimate for the Dunkl heat kernel used in the proof of
%Lemma~\ref{lem:rank-one-decomposition-Dunkl} gives, for every pair of
%multi-indices $\alpha,\beta$,
%\begin{equation}\label{eq:heat-local-Cwikel}
%\left|
%\partial_x^\alpha\partial_y^\beta h_t(x,y)
%\right|
%\lesssim
%t^{-\frac{|\alpha|+|\beta|}{2}}
%V(x,y,\sqrt t)^{-1}
%\exp\left(-c\frac{d(x,y)^2}{t}\right).
%\end{equation}
%Set
%\[
%m:=|\alpha|+|\beta|,
%\qquad
%\delta:=d(x,y)>0.
%\]
%By the integral representation of the fractional power, already justified in
%the proof of Lemma~\ref{lem:rank-one-decomposition-Dunkl},
%\[
%K_{\D}(x,y)
%=
%\frac1{\sqrt\pi}
%\int_0^\infty h_t(x,y)\,\frac{dt}{\sqrt t},
%\]
%and differentiation under the integral sign yields
%\begin{equation}\label{eq:K-derivative-volume-before}
%\left|
%\partial_x^\alpha\partial_y^\beta K_{\D}(x,y)
%\right|
%\lesssim
%\int_0^\infty
%t^{-\frac{m+1}{2}}
%V(x,y,\sqrt t)^{-1}
%\exp\left(-c\frac{\delta^2}{t}\right)\,dt.
%\end{equation}
%
%We estimate the two ranges $0<t\leq\delta^2$ and
%$t\geq\delta^2$ separately.
%
%If $0<t\leq\delta^2$, then $\sqrt t\leq\delta$. By
%\eqref{doublinginequality}, using the upper homogeneous dimension $\N$,
%\[
%V(x,y,\sqrt t)
%\gtrsim
%\left(\frac{\sqrt t}{\delta}\right)^{\N}
%V(x,y,\delta).
%\]
%Consequently,
%\begin{align}
%&\int_0^{\delta^2}
%t^{-\frac{m+1}{2}}
%V(x,y,\sqrt t)^{-1}
%\exp\left(-c\frac{\delta^2}{t}\right)\,dt
%\nonumber\\
%&\qquad\lesssim
%\frac{\delta^\N}{V(x,y,\delta)}
%\int_0^{\delta^2}
%t^{-\frac{\N+m+1}{2}}
%\exp\left(-c\frac{\delta^2}{t}\right)\,dt.
%\label{eq:small-t-Cwikel}
%\end{align}
%Making the change of variables $s=\delta^2/t$, we obtain
%\begin{align*}
%\int_0^{\delta^2}
%t^{-\frac{\N+m+1}{2}}
%e^{-c\delta^2/t}\,dt
%&=
%\delta^{1-\N-m}
%\int_1^\infty
%s^{\frac{\N+m-3}{2}}e^{-cs}\,ds
%\\
%&\lesssim
%\delta^{1-\N-m}.
%\end{align*}
%Thus,
%\begin{equation}\label{eq:small-t-result}
%\int_0^{\delta^2}
%t^{-\frac{m+1}{2}}
%V(x,y,\sqrt t)^{-1}
%e^{-c\delta^2/t}\,dt
%\lesssim
%\frac{\delta^{1-m}}{V(x,y,\delta)}.
%\end{equation}
%
%If $t\geq\delta^2$, then $\sqrt t\geq\delta$. Using the lower
%homogeneous dimension $N$ in \eqref{doublinginequality}, we have
%\[
%V(x,y,\sqrt t)
%\gtrsim
%\left(\frac{\sqrt t}{\delta}\right)^N
%V(x,y,\delta).
%\]
%Hence
%\begin{align}
%&\int_{\delta^2}^{\infty}
%t^{-\frac{m+1}{2}}
%V(x,y,\sqrt t)^{-1}
%e^{-c\delta^2/t}\,dt
%\nonumber\\
%&\qquad\lesssim
%\frac{\delta^N}{V(x,y,\delta)}
%\int_{\delta^2}^{\infty}
%t^{-\frac{N+m+1}{2}}\,dt
%\lesssim
%\frac{\delta^{1-m}}{V(x,y,\delta)}.
%\label{eq:large-t-result}
%\end{align}
%Combining
%\eqref{eq:K-derivative-volume-before},
%\eqref{eq:small-t-result}, and
%\eqref{eq:large-t-result}, we obtain the localized derivative estimate
%\begin{equation}\label{eq:localized-K-derivative}
%\boxed{
%\left|
%\partial_x^\alpha\partial_y^\beta K_{\D}(x,y)
%\right|
%\lesssim
%\frac{d(x,y)^{1-|\alpha|-|\beta|}}
%{V(x,y,d(x,y))}
%}.
%\end{equation}
%
%We now repeat the dyadic localization and Fourier-expansion argument from
%Lemma~\ref{lem:rank-one-decomposition-Dunkl}, but use
%\eqref{eq:localized-K-derivative} instead of the cruder estimate
%\[
%|\partial_x^\alpha\partial_y^\beta K_{\D}(x,y)|
%\lesssim
%d(x,y)^{-(\N-1+|\alpha|+|\beta|)}.
%\]
%
%Fix a Weyl chamber $\Gamma$. Since
%$\operatorname{span}(R)=\mathbb R^N$, Corollary
%\ref{coro:chamber-orthant} gives a linear isomorphism
%\[
%\Phi:\mathbb R^N\longrightarrow\mathbb R^N
%\]
%such that
%\[
%\Phi(\Gamma)=\Omega:=\mathbb R_+^N.
%\]
%Let $\mathbb D_\Gamma$ be the dyadic system constructed in
%Corollary~\ref{Construction over gamma}.
%For
%\[
%I=\Phi^{-1}(Q)\in\mathbb D_\Gamma,
%\]
%write
%\[
%\ell(I)=\ell(Q).
%\]
%
%As in the proof of Lemma~\ref{lem:rank-one-decomposition-Dunkl},
%after fixing $\sigma,\tau\in G$ and localizing by the function $\psi$,
%the relevant pairs $(x,y)$ satisfy
%\begin{equation}\label{eq:distance-scale-Cwikel}
%d(x,y)\simeq\ell(I).
%\end{equation}
%Moreover, all points occurring after introducing the fixed smooth cutoffs in
%the Fourier expansion remain within distance $C\ell(I)$ from
%$\sigma I$ and $\tau I^\ast$, respectively. By the doubling property,
%Theorem~\ref{main:dunkldoubling}, the $G$-invariance of $dw$, and
%\eqref{eq:distance-scale-Cwikel},
%\begin{equation}\label{eq:V-wI-comparable}
%V(x,y,d(x,y))
%\simeq
%w(I)
%\end{equation}
%uniformly on the support of each localized piece. Here, as in
%Lemma~\ref{lem:rank-one-decomposition-Dunkl},
%\[
%I^\ast=\Phi^{-1}((2Q)\cap\Omega),
%\qquad
%w(I^\ast)\simeq w(I).
%\]
%
%Therefore \eqref{eq:localized-K-derivative} implies that, on every such
%localized piece,
%\begin{equation}\label{eq:localized-K-derivative-I}
%\left|
%\partial_x^\alpha\partial_y^\beta K_{\D}(x,y)
%\right|
%\lesssim
%\frac{\ell(I)^{1-|\alpha|-|\beta|}}{w(I)}.
%\end{equation}
%
%Let $c_Q$ be the center of $Q$ and put
%\[
%a=\frac{\Phi(\sigma^{-1}x)-c_Q}{\ell(I)},
%\qquad
%b=\frac{\Phi(\tau^{-1}y)-c_Q}{\ell(I)}.
%\]
%Using exactly the same cutoffs $\chi_Q$ and $\psi$ as in the proof of
%Lemma~\ref{lem:rank-one-decomposition-Dunkl}, define
%\begin{equation}\label{eq:rescaled-H-Cwikel}
%\begin{split}
%\widetilde H_{\sigma,\tau,I}(a,b)
%:={}&
%\frac{w(I)}{\ell(I)}
%\chi_Q(a,b)\psi(a-b)
%\\
%&\times
%K_{\D}\left(
%\sigma\Phi^{-1}(c_Q+\ell(I)a),
%\tau\Phi^{-1}(c_Q+\ell(I)b)
%\right).
%\end{split}
%\end{equation}
%Every derivative of total order $m$ with respect to $(a,b)$ produces a factor
%$\ell(I)^m$. Hence, by
%\eqref{eq:localized-K-derivative-I},
%\[
%\frac{w(I)}{\ell(I)}
%\ell(I)^m
%\frac{\ell(I)^{1-m}}{w(I)}
%\lesssim1.
%\]
%The derivatives of the fixed cutoffs are uniformly bounded. Thus, for every
%integer $M\geq0$,
%\begin{equation}\label{eq:H-uniform-CM-Cwikel}
%\sup_{\sigma,\tau\in G}
%\sup_{I\in\mathbb D_\Gamma}
%\|
%\widetilde H_{\sigma,\tau,I}
%\|_{C^M(\mathbb R^{2N})}
%<\infty.
%\end{equation}
%
%Consequently, the same Fourier-series argument as in
%Lemma~\ref{lem:rank-one-decomposition-Dunkl} gives coefficients
%$b_{p,q,I}^{\sigma,\tau}$ such that
%\[
%\widetilde H_{\sigma,\tau,I}(a,b)
%=
%\sum_{p,q\in\mathbb Z^N}
%b_{p,q,I}^{\sigma,\tau}
%e^{\frac{\pi i}{2}p\cdot a}
%e^{\frac{\pi i}{2}q\cdot b}
%\]
%and, for every $M>0$,
%\begin{equation}\label{eq:bpq-rapid-Cwikel}
%\sup_{\sigma,\tau\in G}
%\sup_{I\in\mathbb D_\Gamma}
%|b_{p,q,I}^{\sigma,\tau}|
%\leq
%C_M(1+|p|+|q|)^{-M}.
%\end{equation}
%
%Fix $M>2N$ and set
%\[
%c_{p,q}:=
%C_M(1+|p|+|q|)^{-M},
%\]
%with $C_M$ chosen large enough that
%\[
%|b_{p,q,I}^{\sigma,\tau}|\leq c_{p,q}
%\]
%for every $\sigma,\tau,I$. Then
%\begin{equation}\label{eq:cpq-l1-Cwikel}
%\sum_{p,q\in\mathbb Z^N}c_{p,q}<\infty.
%\end{equation}
%The quotient
%\[
%a_{p,q,I}^{\sigma,\tau}
%:=
%\frac{b_{p,q,I}^{\sigma,\tau}}{c_{p,q}}
%\]
%satisfies
%\[
%|a_{p,q,I}^{\sigma,\tau}|\leq1.
%\]
%We absorb this factor into the first rank-one function.
%
%For
%\[
%A_{\sigma,I}
%:=
%\|f\|_{L^r(\sigma I,dw)},
%\]
%define, when $A_{\sigma,I}>0$,
%\begin{equation}\label{eq:e-Cwikel}
%\begin{split}
%e_{p,q,I}^{\sigma,\tau}(x)
%:={}&
%a_{p,q,I}^{\sigma,\tau}
%w(I)^{\frac1r-\frac12}
%\frac{f(x)\chi_{\sigma I}(x)}{A_{\sigma,I}}
%\\
%&\times
%\exp\left(
%\frac{\pi i}{2\ell(I)}
%p\cdot
%\bigl(\Phi(\sigma^{-1}x)-c_Q\bigr)
%\right),
%\end{split}
%\end{equation}
%and put $e_{p,q,I}^{\sigma,\tau}=0$ when $A_{\sigma,I}=0$.
%Also define
%\begin{equation}\label{eq:h-Cwikel}
%\begin{split}
%h_{p,q,I}^{\sigma,\tau}(y)
%:={}&
%w(I)^{-1/2}
%\chi_{\tau I^\ast}(y)
%\\
%&\times
%\exp\left(
%-\frac{\pi i}{2\ell(I)}
%q\cdot
%\bigl(\Phi(\tau^{-1}y)-c_Q\bigr)
%\right).
%\end{split}
%\end{equation}
%If necessary, we replace $h_{p,q,I}^{\sigma,\tau}$ by its complex conjugate
%according to the convention
%\[
%(u\otimes v)\xi=\langle\xi,v\rangle u.
%\]
%This has no effect on any support or norm estimate below.
%
%The support properties are
%\begin{equation}\label{eq:support-Cwikel}
%\operatorname{supp}e_{p,q,I}^{\sigma,\tau}
%\subset\sigma I,
%\qquad
%\operatorname{supp}h_{p,q,I}^{\sigma,\tau}
%\subset\tau I^\ast
%\subset\tau(3_\Phi I).
%\end{equation}
%Furthermore,
%\begin{equation}\label{eq:e-Lr-Cwikel}
%\|e_{p,q,I}^{\sigma,\tau}\|_{L^r(dw)}
%\leq
%w(I)^{\frac1r-\frac12},
%\end{equation}
%whereas $w(I^\ast)\simeq w(I)$ implies
%\begin{equation}\label{eq:h-Lr-Cwikel}
%\|h_{p,q,I}^{\sigma,\tau}\|_{L^r(dw)}
%\lesssim
%w(I)^{\frac1r-\frac12}.
%\end{equation}
%
%Define
%\begin{equation}\label{eq:lambda-Cwikel}
%\lambda_{\sigma,I}(f)
%:=
%\ell(I)\,
%w(I)^{-1/r}
%\|f\|_{L^r(\sigma I,dw)}.
%\end{equation}
%A direct calculation gives
%\[
%\lambda_{\sigma,I}(f)
%w(I)^{\frac1r-\frac12}
%A_{\sigma,I}^{-1}
%w(I)^{-1/2}
%=
%\frac{\ell(I)}{w(I)}.
%\]
%It follows from the Fourier expansion above and the dyadic partition of unity
%used in Lemma~\ref{lem:rank-one-decomposition-Dunkl} that
%\begin{equation}\label{eq:operator-rank-one-Cwikel}
%M_f(-\Delta_{\D})^{-1/2}
%=
%C
%\sum_{\sigma,\tau\in G}
%\sum_{p,q\in\mathbb Z^N}
%c_{p,q}
%\sum_{I\in\mathbb D_\Gamma}
%\lambda_{\sigma,I}(f)\,
%e_{p,q,I}^{\sigma,\tau}
%\otimes
%h_{p,q,I}^{\sigma,\tau},
%\end{equation}
%initially as an identity of kernels, and hence as an identity of bilinear forms
%on the natural test-function class.
%
%This is the localized refinement of
%Lemma~\ref{lem:rank-one-decomposition-Dunkl} needed below. Notice that
%the power $\ell(I)^{-(\N-1)}$ from the cruder version has disappeared:
%the local size of the fractional kernel is
%\[
%\frac{\ell(I)}{w(I)},
%\]
%which is the correct scale for the present Cwikel estimate.
%
%\medskip
%\noindent
%\textbf{Step 2. A variant of XXX.}
%
%We prove
%\begin{equation}\label{eq:lambda-weak-lN}
%\left\|
%\left\{
%\lambda_{\sigma,I}(f)
%\right\}_{I\in\mathbb D_\Gamma}
%\right\|_{\ell^{N,\infty}}
%\lesssim
%\|f\|_{L^N(\mathbb R^N,dx)}
%\end{equation}
%uniformly in $\sigma\in G$.
%
%Since $\operatorname{span}(R)=\mathbb R^N$, we have $N_0=N$ in
%Lemma~\ref{mulemma}. Thus
%\[
%\Omega=\mathbb R_+^N,
%\]
%and the push-forward measure
%\[
%\mu(E):=w(\Phi^{-1}(E))
%\]
%has density
%\begin{equation}\label{eq:phi-Cwikel}
%d\mu(u)=\phi(u)\,du,
%\qquad
%\phi(u)
%=
%|\det\Phi|^{-1}
%\prod_{v\in R_+}
%\left(
%\sum_{j=1}^N a_{v,j}u_j
%\right)^{2k(v)},
%\end{equation}
%where
%\[
%a_{v,j}\geq0.
%\]
%Hence $\phi$ is nondecreasing in each coordinate.
%
%Let
%\[
%Q
%=
%\prod_{j=1}^N
%[m_j\ell,(m_j+1)\ell)
%\subset\Omega,
%\qquad m_j\in\mathbb Z_+,
%\]
%and let
%\[
%u^+
%:=
%((m_1+1)\ell,\ldots,(m_N+1)\ell)
%\]
%be its upper corner. Set
%\[
%Q^+
%:=
%\prod_{j=1}^N
%[(m_j+\tfrac12)\ell,(m_j+1)\ell).
%\]
%Then
%\[
%|Q^+|=2^{-N}|Q|.
%\]
%For $u\in Q^+$,
%\[
%u_j
%\geq
%(m_j+\tfrac12)\ell
%\geq
%\frac12(m_j+1)\ell
%=
%\frac12 u_j^+.
%\]
%Since $a_{v,j}\geq0$,
%\[
%\sum_{j=1}^N a_{v,j}u_j
%\geq
%\frac12
%\sum_{j=1}^N a_{v,j}u_j^+.
%\]
%Therefore
%\[
%\phi(u)
%\geq
%2^{-2\sum_{v\in R_+}k(v)}
%\phi(u^+).
%\]
%Since
%\[
%2\sum_{v\in R_+}k(v)
%=
%\sum_{v\in R}k(v)
%=
%\N-N,
%\]
%we obtain
%\[
%\phi(u)
%\geq
%2^{-(\N-N)}\phi(u^+),
%\qquad u\in Q^+.
%\]
%Consequently,
%\begin{align}
%\mu(Q)
%&=
%\int_Q\phi(u)\,du
%\nonumber\\
%&\geq
%\int_{Q^+}\phi(u)\,du
%\nonumber\\
%&\geq
%2^{-(\N-N)}
%|Q^+|\phi(u^+)
%\nonumber\\
%&=
%2^{-\N}|Q|\phi(u^+).
%\label{eq:phi-average-lower}
%\end{align}
%Since $\phi$ is coordinatewise nondecreasing,
%\[
%\sup_{u\in Q}\phi(u)=\phi(u^+)
%\]
%in the sense of essential supremum. Hence
%\begin{equation}\label{eq:phi-sup-average}
%\boxed{
%\sup_{u\in Q}\phi(u)
%\leq
%2^\N\frac{\mu(Q)}{|Q|}
%}.
%\end{equation}
%
%Let $I=\Phi^{-1}(Q)$. By the $G$-invariance of $dw$,
%\begin{align*}
%\frac1{w(I)}
%\int_{\sigma I}|f(x)|^r\,dw(x)
%&=
%\frac1{\mu(Q)}
%\int_Q
%|f(\sigma\Phi^{-1}(u))|^r
%\phi(u)\,du
%\\
%&\leq
%\frac{\sup_Q\phi}{\mu(Q)}
%\int_Q
%|f(\sigma\Phi^{-1}(u))|^r\,du
%\\
%&\lesssim
%\frac1{|Q|}
%\int_Q
%|f(\sigma\Phi^{-1}(u))|^r\,du,
%\end{align*}
%where the last inequality follows from
%\eqref{eq:phi-sup-average}. Therefore,
%by \eqref{eq:lambda-Cwikel},
%\begin{align}
%\lambda_{\sigma,I}(f)^r
%&=
%\ell(I)^r
%\frac1{w(I)}
%\int_{\sigma I}|f(x)|^r\,dw(x)
%\nonumber\\
%&\lesssim
%\ell(Q)^r
%\frac1{|Q|}
%\int_Q
%|f(\sigma\Phi^{-1}(u))|^r\,du.
%\label{eq:lambda-to-Lebesgue-average}
%\end{align}
%Since $|Q|=\ell(Q)^N$,
%\[
%\ell(Q)^r|Q|^{-1}
%=
%|Q|^{\frac rN-1}.
%\]
%Set
%\[
%s:=\frac Nr.
%\]
%By \eqref{eq:r-range-Cwikel},
%\[
%1<s<\infty.
%\]
%Define on $\mathbb R^N$
%\[
%F_\sigma(u)
%:=
%\begin{cases}
%|f(\sigma\Phi^{-1}(u))|^r,
%&u\in\Omega,\\
%0,
%&u\notin\Omega.
%\end{cases}
%\]
%Then \eqref{eq:lambda-to-Lebesgue-average} becomes
%\begin{equation}\label{eq:lambda-r-Zhijie}
%\lambda_{\sigma,I}(f)^r
%\lesssim
%|Q|^{\frac1s-1}
%\int_QF_\sigma(u)\,du.
%\end{equation}
%
%The family $\mathbb D_\Omega^0$ is a subfamily of the standard dyadic grid
%in $\mathbb R^N$. Hence Lemma XXX, applied to $F_\sigma$
%with exponent $s=N/r$, gives
%\begin{align}
%\left\|
%\left\{
%\lambda_{\sigma,I}(f)^r
%\right\}_{I\in\mathbb D_\Gamma}
%\right\|_{\ell^{N/r,\infty}}
%&\lesssim
%\|F_\sigma\|_{L^{N/r}(\mathbb R^N)}.
%\label{eq:Zhijie-applied}
%\end{align}
%
%For every nonnegative sequence $\{a_I\}$,
%\begin{equation}\label{eq:weak-power-identity}
%\left\|
%\{a_I^r\}
%\right\|_{\ell^{N/r,\infty}}
%=
%\left\|
%\{a_I\}
%\right\|_{\ell^{N,\infty}}^r.
%\end{equation}
%Indeed,
%\[
%\#\{I:a_I^r>t\}
%=
%\#\{I:a_I>t^{1/r}\},
%\]
%and the identity follows directly from the distribution-function definition of
%the weak Lorentz norm.
%
%Furthermore,
%\begin{align*}
%\|F_\sigma\|_{L^{N/r}}^{1/r}
%&=
%\left(
%\int_\Omega
%|f(\sigma\Phi^{-1}(u))|^N\,du
%\right)^{1/N}
%\\
%&=
%|\det\Phi|^{1/N}
%\left(
%\int_{\sigma\Gamma}|f(x)|^N\,dx
%\right)^{1/N}
%\\
%&\leq
%|\det\Phi|^{1/N}
%\|f\|_{L^N(\mathbb R^N,dx)}.
%\end{align*}
%Combining this with
%\eqref{eq:Zhijie-applied} and
%\eqref{eq:weak-power-identity} proves
%\eqref{eq:lambda-weak-lN}.
%
%\medskip
%\noindent
%\textbf{Step 3. Nearly weak orthogonality of the rank-one factors.}
%
%Fix
%\[
%\sigma,\tau\in G,
%\qquad
%p,q\in\mathbb Z^N.
%\]
%Let
%\[
%U_\sigma:L^2(\Gamma,dw)\longrightarrow L^2(\sigma\Gamma,dw),
%\qquad
%U_\sigma u(x):=u(\sigma^{-1}x).
%\]
%Since $dw$ is $G$-invariant, $U_\sigma$ is unitary. Define the pullbacks
%\[
%\widetilde e_I
%:=
%U_\sigma^{-1}
%e_{p,q,I}^{\sigma,\tau},
%\qquad
%\widetilde h_I
%:=
%U_\tau^{-1}
%h_{p,q,I}^{\sigma,\tau}.
%\]
%Then
%\[
%\operatorname{supp}\widetilde e_I\subset I,
%\qquad
%\operatorname{supp}\widetilde h_I\subset I^\ast.
%\]
%By the construction of $\mathbb D_\Gamma$, the bi-Lipschitz property
%\eqref{eq:biLip-N}, and $I^\ast\subset3_\Phi I$, there is a fixed constant
%$C_\ast$ such that
%\[
%I\cup I^\ast
%\subset
%B_I^\ast
%:=
%B_\Gamma(z_I,C_\ast\ell(I))
%\]
%for a suitable $z_I\in I$. Moreover, by
%Theorem~\ref{main:dunkldoubling},
%\begin{equation}\label{eq:Bstar-I-measure}
%w(B_I^\ast)\simeq w(I).
%\end{equation}
%From \eqref{eq:e-Lr-Cwikel} and
%\eqref{eq:h-Lr-Cwikel},
%\begin{equation}\label{eq:NWO-Lr}
%\|\widetilde e_I\|_{L^r(dw)}
%+
%\|\widetilde h_I\|_{L^r(dw)}
%\lesssim
%w(I)^{\frac1r-\frac12}.
%\end{equation}
%
%For $u\in L^2(\Gamma,dw)$, define
%\[
%\mathcal M_1u(x)
%:=
%\sup_{I\in\mathbb D_\Gamma}
%\chi_I(x)
%\frac{|\langle u,\widetilde e_I\rangle|}
%{w(I)^{1/2}}.
%\]
%By H\"older's inequality, \eqref{eq:NWO-Lr}, and
%\eqref{eq:Bstar-I-measure},
%\begin{align*}
%\frac{|\langle u,\widetilde e_I\rangle|}{w(I)^{1/2}}
%&\leq
%\frac{
%\|\chi_{B_I^\ast}u\|_{L^{r'}(dw)}
%\|\widetilde e_I\|_{L^r(dw)}
%}{w(I)^{1/2}}
%\\
%&\lesssim
%\|\chi_{B_I^\ast}u\|_{L^{r'}(dw)}
%w(I)^{\frac1r-1}
%\\
%&=
%\|\chi_{B_I^\ast}u\|_{L^{r'}(dw)}
%w(I)^{-1/r'}
%\\
%&\lesssim
%\left(
%\fint_{B_I^\ast}|u|^{r'}\,dw
%\right)^{1/r'}.
%\end{align*}
%Therefore
%\begin{equation}\label{eq:M1-maximal}
%\mathcal M_1u
%\lesssim
%\left(
%M_\Gamma(|u|^{r'})
%\right)^{1/r'},
%\end{equation}
%where $M_\Gamma$ denotes the Hardy--Littlewood maximal operator on the
%doubling metric measure space $(\Gamma,\|\cdot\|,dw)$.
%
%Since $r'<2$, we have
%\[
%\frac2{r'}>1.
%\]
%The Hardy--Littlewood maximal operator is bounded on
%$L^{2/r'}(\Gamma,dw)$, and hence
%\begin{align*}
%\|\mathcal M_1u\|_{L^2(dw)}
%&\lesssim
%\left\|
%M_\Gamma(|u|^{r'})
%\right\|_{L^{2/r'}(dw)}^{1/r'}
%\\
%&\lesssim
%\||u|^{r'}\|_{L^{2/r'}(dw)}^{1/r'}
%\\
%&=
%\|u\|_{L^2(dw)}.
%\end{align*}
%The identical argument, with $\widetilde h_I$ in place of
%$\widetilde e_I$, proves the $L^2$ boundedness of
%\[
%\mathcal M_2v(x)
%:=
%\sup_{I\in\mathbb D_\Gamma}
%\chi_I(x)
%\frac{|\langle v,\widetilde h_I\rangle|}
%{w(I)^{1/2}}.
%\]
%
%Now define the bi-sublinear maximal operator
%\[
%\mathcal M_{\mathcal E}(u,v)(x)
%:=
%\sup_{I\in\mathbb D_\Gamma}
%\chi_I(x)
%\frac{
%|\langle u,\widetilde e_I\rangle
%\langle v,\widetilde h_I\rangle|
%}{w(I)}.
%\]
%Then
%\[
%\mathcal M_{\mathcal E}(u,v)
%\leq
%\mathcal M_1u\,\mathcal M_2v.
%\]
%Thus, by Cauchy--Schwarz,
%\begin{equation}\label{eq:ME-L2L2-L1}
%\|
%\mathcal M_{\mathcal E}(u,v)
%\|_{L^1(dw)}
%\lesssim
%\|u\|_{L^2(dw)}
%\|v\|_{L^2(dw)}.
%\end{equation}
%This is precisely the NWO estimate required in
%\cite[Corollary~7.6 and Remark~7.9]{HKarxiv}.
%
%\medskip
%\noindent
%\textbf{Step 4. Application of Corollary 7.6.}
%
%For the fixed $\sigma,\tau,p,q$, consider on
%$L^2(\Gamma,dw)$ the operator
%\[
%A_{\sigma,\tau,p,q}(f)
%:=
%\sum_{I\in\mathbb D_\Gamma}
%\lambda_{\sigma,I}(f)
%\widetilde e_I\otimes\widetilde h_I.
%\]
%Applying \cite[Corollary~7.6]{HKarxiv} with
%\[
%p=N,
%\qquad
%q=\infty,
%\]
%and using \eqref{eq:ME-L2L2-L1}, we obtain
%\begin{align}
%\|
%A_{\sigma,\tau,p,q}(f)
%\|_{S^{N,\infty}(L^2(\Gamma,dw))}
%&\lesssim
%\left\|
%\{\lambda_{\sigma,I}(f)\}_{I\in\mathbb D_\Gamma}
%\right\|_{\ell^{N,\infty}}
%\nonumber\\
%&\lesssim
%\|f\|_{L^N(\mathbb R^N,dx)},
%\label{eq:block-SNweak}
%\end{align}
%where the last inequality follows from
%\eqref{eq:lambda-weak-lN}.
%
%The original block
%\[
%\sum_{I\in\mathbb D_\Gamma}
%\lambda_{\sigma,I}(f)
%e_{p,q,I}^{\sigma,\tau}
%\otimes
%h_{p,q,I}^{\sigma,\tau}
%\]
%acts from $L^2(\tau\Gamma,dw)$ to
%$L^2(\sigma\Gamma,dw)$ and is obtained from
%$A_{\sigma,\tau,p,q}(f)$ by unitary transformations on the
%domain and the range. Hence its nonzero singular values coincide with those
%of $A_{\sigma,\tau,p,q}(f)$. Therefore
%\begin{equation}\label{eq:original-block-SNweak}
%\left\|
%\sum_{I\in\mathbb D_\Gamma}
%\lambda_{\sigma,I}(f)
%e_{p,q,I}^{\sigma,\tau}
%\otimes
%h_{p,q,I}^{\sigma,\tau}
%\right\|_{S^{N,\infty}(L^2(\mathbb R^N,dw))}
%\lesssim
%\|f\|_{L^N(dx)}.
%\end{equation}
%
%\medskip
%\noindent
%\textbf{Step 5. Summation over the Fourier indices and Weyl chambers.}
%
%For completeness, we recall a convenient norm equivalent to the usual weak
%Schatten quasi-norm. If $T$ is compact, let
%\[
%s_1(T)\geq s_2(T)\geq\cdots
%\]
%be its singular values and define
%\[
%\|T\|_{(N,\infty)}
%:=
%\sup_{m\geq1}
%m^{\frac1N-1}
%\sum_{j=1}^m s_j(T).
%\]
%By the Ky Fan inequality,
%\[
%\sum_{j=1}^m s_j(A+B)
%\leq
%\sum_{j=1}^m s_j(A)
%+
%\sum_{j=1}^m s_j(B),
%\]
%so $\|\cdot\|_{(N,\infty)}$ satisfies the triangle inequality.
%
%Moreover,
%\[
%m^{1/N}s_m(T)
%\leq
%m^{\frac1N-1}
%\sum_{j=1}^m s_j(T),
%\]
%while, if
%\[
%K:=\sup_{j\geq1}j^{1/N}s_j(T),
%\]
%then
%\[
%s_j(T)\leq K j^{-1/N},
%\]
%and hence
%\[
%m^{\frac1N-1}
%\sum_{j=1}^m s_j(T)
%\leq
%K
%m^{\frac1N-1}
%\sum_{j=1}^m j^{-1/N}
%\lesssim_N K.
%\]
%Thus
%\begin{equation}\label{eq:equiv-weak-Schatten-norm}
%\|T\|_{(N,\infty)}
%\simeq_N
%\|T\|_{S^{N,\infty}}.
%\end{equation}
%
%Using the decomposition
%\eqref{eq:operator-rank-one-Cwikel},
%the triangle inequality for $\|\cdot\|_{(N,\infty)}$,
%\eqref{eq:original-block-SNweak}, and
%\eqref{eq:cpq-l1-Cwikel}, we obtain
%\begin{align*}
%&
%\|
%M_f(-\Delta_{\D})^{-1/2}
%\|_{S^{N,\infty}(L^2(dw))}
%\\
%&\qquad\lesssim_N
%\sum_{\sigma,\tau\in G}
%\sum_{p,q\in\mathbb Z^N}
%c_{p,q}
%\left\|
%\sum_{I\in\mathbb D_\Gamma}
%\lambda_{\sigma,I}(f)
%e_{p,q,I}^{\sigma,\tau}
%\otimes
%h_{p,q,I}^{\sigma,\tau}
%\right\|_{S^{N,\infty}}
%\\
%&\qquad\lesssim
%\sum_{\sigma,\tau\in G}
%\sum_{p,q\in\mathbb Z^N}
%c_{p,q}
%\|f\|_{L^N(dx)}
%\\
%&\qquad=
%|G|^2
%\left(
%\sum_{p,q\in\mathbb Z^N}c_{p,q}
%\right)
%\|f\|_{L^N(dx)}
%\\
%&\qquad\lesssim_{N,R,k}
%\|f\|_{L^N(\mathbb R^N,dx)}.
%\end{align*}
%This proves \eqref{eq:Cwikel-Dunkl} for, say, bounded compactly supported
%$f$.
%
%Finally, let $f\in L^N(\mathbb R^N,dx)$ be arbitrary and choose bounded
%compactly supported functions $f_j$ such that
%\[
%\|f_j-f\|_{L^N(dx)}\longrightarrow0.
%\]
%Applying the estimate already proved to $f_j-f_m$, we see that
%\[
%\left\{
%M_{f_j}(-\Delta_{\D})^{-1/2}
%\right\}_{j\geq1}
%\]
%is Cauchy in $S^{N,\infty}$ with respect to the equivalent norm
%\eqref{eq:equiv-weak-Schatten-norm}. It therefore converges to an operator
%$T_f\in S^{N,\infty}$ satisfying
%\[
%\|T_f\|_{S^{N,\infty}}
%\lesssim_{N,R,k}
%\|f\|_{L^N(dx)}.
%\]
%Since $r<N$, every $L^N(dx)$ function belongs locally to $L^r(dw)$:
%indeed, the density $dw/dx$ is bounded on compact subsets of
%$\mathbb R^N$, and $L^N$ embeds locally into $L^r$. Consequently the
%kernel decomposition above is valid locally for $f$. Passing to the limit in
%the kernel identity on the natural test-function class shows that $T_f$ is the
%closure of the integral operator with kernel
%\[
%f(x)K_{\D}(x,y),
%\]
%that is,
%\[
%T_f=M_f(-\Delta_{\D})^{-1/2}.
%\]
%The proof is complete.
%\end{proof}

\section{Cwikel's estimate in $N>2$}

\begin{lemma}\label{lem:Dunkl-fractional-kernel-derivative}
Assume that $N>1$. Then for every pair of multi-indices $\alpha,\beta$, there exists
a constant $C_{\alpha,\beta}>0$ such that
\begin{equation}\label{eq:localized-K-derivative}
\left|
\partial_x^\alpha\partial_y^\beta
K_{(-\Delta_{\D})^{-1/2}}(x,y)
\right|
\leq
C_{\alpha,\beta}
\frac{
d(x,y)^{1-|\alpha|-|\beta|}
}{
V(x,y,d(x,y))
}
\end{equation}
for all $x,y\in\mathbb R^N$ satisfying $d(x,y)>0$.
\end{lemma}

\begin{proof}
Let $h_t(x,y)$ denote the integral kernel of the Dunkl heat semigroup
$e^{t\Delta_{\D}}$. By
\cite[Theorem~4.1(d)]{MR4014802}, for every pair of multi-indices
$\alpha,\beta$, there exist constants
$C_{\alpha,\beta},c_{\alpha,\beta}>0$ such that
\begin{equation}\label{eq:Dunkl-heat-mixed-localized}
\left|
\partial_x^\alpha\partial_y^\beta h_t(x,y)
\right|
\leq
C_{\alpha,\beta}
t^{-\frac{|\alpha|+|\beta|}{2}}
V(x,y,\sqrt t)^{-1}
\exp\left(
-c_{\alpha,\beta}\frac{d(x,y)^2}{t}
\right)
\end{equation}
for all $t>0$ and $x,y\in\mathbb R^N$.

%By \eqref{twosideinequ},
%for every $z\in\mathbb R^N$ and $s>0$,
%\[
%w(B(z,s))
%\gtrsim
%s^N
%\prod_{v\in R}
%\bigl(|\langle z,v\rangle|+s\bigr)^{k(v)}
%\geq
%s^N\prod_{v\in R}s^{k(v)}
%=
%s^{\N}.
%\]
%Hence
%\begin{equation}\label{eq:V-global-lower-fractional}
%V(x,y,\sqrt t)
%\gtrsim
%t^{\N/2},
%\qquad t>0,
%\end{equation}
%uniformly in $x,y\in\mathbb R^N$.

By the spectral functional calculus,
\[
(-\Delta_{\D})^{-1/2}
=
\frac{1}{\sqrt{\pi}}
\int_0^\infty
e^{t\Delta_{\D}}\,
\frac{dt}{\sqrt t}.
\]
Consequently, for $d(x,y)>0$,
\begin{equation}\label{eq:Dunkl-fractional-kernel-localized}
K_{(-\Delta_{\D})^{-1/2}}(x,y)
=
\frac{1}{\sqrt{\pi}}
\int_0^\infty
h_t(x,y)\,\frac{dt}{\sqrt t}.
\end{equation}

It follows that
$K_{(-\Delta_{\D})^{-1/2}}$ is smooth whenever $d(x,y)>0$ and
\begin{equation}\label{eq:Dunkl-fractional-derivative-integral}
\partial_x^\alpha\partial_y^\beta
K_{(-\Delta_{\D})^{-1/2}}(x,y)
=
\frac{1}{\sqrt{\pi}}
\int_0^\infty
\partial_x^\alpha\partial_y^\beta h_t(x,y)
\,\frac{dt}{\sqrt t}.
\end{equation}
Set
$
m:=|\alpha|+|\beta|.
$
Fix $x,y\in\mathbb R^N$ with $d(x,y)>0$. By
\eqref{eq:Dunkl-heat-mixed-localized} and
\eqref{eq:Dunkl-fractional-derivative-integral},
\begin{equation}\label{eq:K-derivative-volume}
\left|
\partial_x^\alpha\partial_y^\beta
K_{(-\Delta_{\D})^{-1/2}}(x,y)
\right|
\lesssim
\int_0^\infty
t^{-\frac{m+1}{2}}
V(x,y,\sqrt t)^{-1}
\exp\left(
-c\frac{d(x,y)^2}{t}
\right)\,dt.
\end{equation}

Suppose first that
$
0<t\leq d(x,y)^2.
$
Then $\sqrt t\leq d(x,y)$. By the upper estimate in
\eqref{doublinginequality}, for $z=x,y$,
\[
\frac{w(B(z,d(x,y)))}
{w(B(z,\sqrt t))}
\lesssim
\left(
\frac{d(x,y)}{\sqrt t}
\right)^\N.
\]
Equivalently,
\[
w(B(z,\sqrt t))
\gtrsim
\left(
\frac{\sqrt t}{d(x,y)}
\right)^\N
w(B(z,d(x,y))).
\]
Taking the maximum over $z=x,y$ gives
\begin{equation}\label{eq:V-small-scale-comparison}
V(x,y,\sqrt t)
\gtrsim
\left(
\frac{\sqrt t}{d(x,y)}
\right)^\N
V(x,y,d(x,y)).
\end{equation}
Therefore,
\begin{align*}
&\int_0^{d(x,y)^2}
t^{-\frac{m+1}{2}}
V(x,y,\sqrt t)^{-1}
\exp\left(
-c\frac{d(x,y)^2}{t}
\right)\,dt
\\
&\qquad\lesssim
\frac{d(x,y)^\N}{V(x,y,d(x,y))}
\int_0^{d(x,y)^2}
t^{-\frac{\N+m+1}{2}}
\exp\left(
-c\frac{d(x,y)^2}{t}
\right)\,dt.
\end{align*}
Using again the change of variables
$
s=\frac{d(x,y)^2}{t},
$
we obtain
\begin{align*}
&\int_0^{d(x,y)^2}
t^{-\frac{\N+m+1}{2}}
\exp\left(
-c\frac{d(x,y)^2}{t}
\right)\,dt=
d(x,y)^{1-\N-m}
\int_1^\infty
s^{\frac{\N+m-3}{2}}e^{-cs}\,ds\lesssim
d(x,y)^{1-\N-m}.
\end{align*}
Consequently,
\begin{equation}\label{eq:K-small-time-localized}
\int_0^{d(x,y)^2}
t^{-\frac{m+1}{2}}
V(x,y,\sqrt t)^{-1}
\exp\left(
-c\frac{d(x,y)^2}{t}
\right)\,dt
\lesssim
\frac{d(x,y)^{1-m}}{V(x,y,d(x,y))}.
\end{equation}

Suppose next that
$
t\geq d(x,y)^2.
$
Then $\sqrt t\geq d(x,y)$. By the lower estimate in
\eqref{doublinginequality}, for $z=x,y$,
\[
\frac{w(B(z,\sqrt t))}
{w(B(z,d(x,y)))}
\gtrsim
\left(
\frac{\sqrt t}{d(x,y)}
\right)^N.
\]
It follows that
\begin{equation}\label{eq:V-large-scale-comparison}
V(x,y,\sqrt t)
\gtrsim
\left(
\frac{\sqrt t}{d(x,y)}
\right)^N
V(x,y,d(x,y)).
\end{equation}
Hence
\begin{align*}
&\int_{d(x,y)^2}^{\infty}
t^{-\frac{m+1}{2}}
V(x,y,\sqrt t)^{-1}
\exp\left(
-c\frac{d(x,y)^2}{t}
\right)\,dt\lesssim
\frac{d(x,y)^N}{V(x,y,d(x,y))}
\int_{d(x,y)^2}^{\infty}
t^{-\frac{N+m+1}{2}}\,dt.
\end{align*}
Since $N>1$ and $m\geq0$, we have $N+m>1$, and therefore
\[
\int_{d(x,y)^2}^{\infty}
t^{-\frac{N+m+1}{2}}\,dt
=
\frac{2}{N+m-1}
d(x,y)^{1-N-m}.
\]
Thus
\begin{equation}\label{eq:K-large-time-localized}
\int_{d(x,y)^2}^{\infty}
t^{-\frac{m+1}{2}}
V(x,y,\sqrt t)^{-1}
\exp\left(
-c\frac{d(x,y)^2}{t}
\right)\,dt
\lesssim
\frac{d(x,y)^{1-m}}{V(x,y,d(x,y))}.
\end{equation}

Combining
\eqref{eq:K-derivative-volume},
\eqref{eq:K-small-time-localized}, and
\eqref{eq:K-large-time-localized}, we obtain
\[
\left|
\partial_x^\alpha\partial_y^\beta
K_{(-\Delta_{\D})^{-1/2}}(x,y)
\right|
\lesssim
\frac{d(x,y)^{1-m}}{V(x,y,d(x,y))}.
\]
Since $m=|\alpha|+|\beta|$, this gives
\[
\left|
\partial_x^\alpha\partial_y^\beta
K_{(-\Delta_{\D})^{-1/2}}(x,y)
\right|
\lesssim
\frac{
d(x,y)^{1-|\alpha|-|\beta|}
}{
V(x,y,d(x,y))
},
\]
which is exactly \eqref{eq:localized-K-derivative}.
\end{proof}

\begin{lemma}\label{lem:Cwikel-Dunkl}
Assume that $N>2.$ There exists a constant $C_{N,R,k}>0$ such that, for every
$f\in L_N(\mathbb{R}^N,dx),$
$$M_f(-\Delta_{\D})^{-\frac12}\in
\mathcal{L}_{N,\infty}\bigl(L_2(\mathbb{R}^N,w(x)dx)\bigr)$$
and
\begin{equation}\label{eq:Cwikel-Dunkl}
\bigl\|M_f(-\Delta_{\D})^{-\frac12}\bigr\|_
{\mathcal{L}_{N,\infty}(L_2(\mathbb{R}^N,w(x)dx))}\leq C_{N,R,k}\|f\|_{L_N(\mathbb{R}^N,dx)}.
\end{equation}
\end{lemma}
\begin{proof} For $g\in G,$ let $U_g:\xi\to\xi\circ g^{-1}.$ For a chamber $\Gamma$ and $g\in G,$ define an integral operator
$$V_{\Gamma,g}=M_{\chi_{\Gamma}}(-\Delta_{\D})^{-\frac12}M_{\chi_{g\Gamma}}U_g.$$
By Lemma \ref{lem:Dunkl-fractional-kernel-derivative}, each operator $V_{\Gamma,g}$ satisfies the assumptions in Theorem \ref{simple cwikel estimate}. We, obviously, have
$$M_f(-\Delta_{\D})^{-\frac12}=\sum_{\Gamma}\sum_{g\in G}M_fV_{\Gamma,g}\cdot U_{g^{-1}}.$$
The assertion follows now from Theorem \ref{simple cwikel estimate}.
\end{proof}


\section{Better Cwikel estimate}

\begin{theorem}\label{simple cwikel estimate} Let $\Gamma$ be a chamber. Let $V$ be an integral operator on $L_2(\Gamma,w(x)dx)$ with integral kernel $K$ satisfying the conditions
$$|\partial_x^{\alpha}\partial_y^{\beta}K(x,y)|\leq C_{\alpha,\beta}\frac{|x-y|^{1-|\alpha|_1-|\beta|_1}}{V(x,|x-y|)}.$$
We have
$$\|M_fV\|_{\mathcal{L}_{N,\infty}(\Gamma,w(x)dx)}\leq c_K\|f\|_{L_N(\Gamma)}.$$
\end{theorem}

\begin{nota}
Fix a closed Weyl chamber $\Gamma$ and let $\Phi:\mathbb{R}^N\to\mathbb{R}^N$ be the linear isomorphism appearing in
Corollary \ref{coro:chamber-orthant}.


Let $\mathbb{D}_{\Gamma}$ be the dyadic system over $\Gamma$ constructed in
Corollary \ref{Construction over gamma}.
If $I\in\mathbb{D}_{\Gamma}$ and if $Q=\Phi(I),$ then we denote
$$d_I=\ell(Q).$$


For every $I\in\mathbb{D}_{\Gamma}$ and for every $f\in L_r^{\mathrm{loc}}(\mathbb{R} ^N,w(x)dx)$, define
\begin{equation}\label{eq:Dunkl-lambda-rank-one}
\lambda_I(f)=d_Iw(I)^{-\frac1r}\|f\|_{L_r(I,w(x)dx)}.
\end{equation}
\end{nota}

\begin{lemma} There exists a smooth function $\psi$ supported on $[-1,1]^N\backslash(-\frac14,\frac14)^N$ such that
$$1=\sum_{I\in\mathbb{D}_{\Gamma}}\chi_I(x)\chi_{3I}(y)\psi(\frac{\Phi(x)-\Phi(y)}{d_I}).$$
\end{lemma}
\begin{proof} Choose $\rho\in C_c^{\infty}(\mathbb{R}^N)$ such that $\rho=1$ on $(-\frac12,\frac12)^N$ and $\rho=0$ outside $(-1,1)^N.$ Set $\psi(z)=\rho(z)-\rho(2z).$ Clearly,
\begin{equation}\label{eq:rankone-psi-support}
\operatorname{supp}(\psi)\subset
\left\{z\in\mathbb{R}^N:\quad 
\frac14\leq |z|_{\infty}\leq 1\right\}.
\end{equation}
For every $z\neq0,$ we have
\begin{equation}\label{eq:rankone-psi-partition}
\sum_{j\in\mathbb{Z}}\psi(2^jz)=1.
\end{equation}

For every $x\in\Gamma,$ there exists a unique $I\in\mathbb{D}_{\Gamma,j}$ containing $x.$ For this $I,$ we have $d_I=2^{-j}.$ We have
$$1=\sum_{j\in\mathbb{Z}}\psi(2^j(\Phi(x)-\Phi(y)))=\sum_{j\in\mathbb{Z}}(\sum_{I\in\mathbb{D}_{\Gamma,j}}\chi_I(x))\psi(2^j(\Phi(x)-\Phi(y)))=$$
$$=\sum_{j\in\mathbb{Z}}\sum_{I\in\mathbb{D}_{\Gamma,j}}\chi_I(x)\psi(\frac{\Phi(x)-\Phi(y)}{d_I})=\sum_{I\in\mathbb{D}_{\Gamma}}\chi_I(x)\psi(\frac{\Phi(x)-\Phi(y)}{d_I}).$$

If $x\in I$ and $\psi(\frac{\Phi(x)-\Phi(y)}{d_I})\neq 0,$ then
$$\Phi(x)\in\Phi(I)=Q,\quad |\Phi(x)-\Phi(y)|\leq d_I=\ell(Q).$$
We obviously have $|\Phi(x)-c_Q|_{\infty}\leq\frac12\ell(Q)$ and $|\Phi(x)-\Phi(y)|_{\infty}\leq\ell(Q).$ Thus, $|\Phi(y)-c_Q|\leq\frac32\ell(Q).$ In other words, $\Phi(y)\in 3Q$ or, equivalently, $y\in 3I.$ Thus,
$$\chi_I(x)\psi(\frac{\Phi(x)-\Phi(y)}{d_I})=\chi_I(x)\chi_{3I}(y)\psi(\frac{\Phi(x)-\Phi(y)}{d_I}).$$
This completes the proof.
\end{proof}


\begin{lemma}\label{lem:rank-one-decomposition-Dunkl} Assume that $N>1$ and $2<r<\infty.$ Let $f\in L_r^{\mathrm{loc}}(\mathbb{R} ^N,w(x)dx).$ There exist functions
$
\{e_{k,l,I}\}_{k,l\in\mathbb{Z}^N,I\in\mathbb{D}_{\Gamma}},$ $
\{f_{k,l,I}\}_{k,l\in\mathbb{Z}^N,I\in\mathbb{D}_{\Gamma}},$
such that
\begin{enumerate}
\item $e_{k,l,I}$ is supported in $I$ and $f_{k,l,I}$ is supported in $3I;$
\item
$$\|e_{k,l,I}\|_{L_r(\mathbb{R}^N,w(x)dx)}=\|f\|_{L_r(\mathbb{R}^N,w(x)dx)},\quad \|f_{k,l,I}\|_{L_r(\mathbb{R}^N,w(x)dx)}\leq c_{K,R,k}w(I)^{\frac1r};$$
\item For almost every $x,y\in\Gamma,$ we have
\begin{equation}\label{eq:Dunkl-kernel-rank-one-new}
f(x)K(x,y)=\sum_{k,l\in\mathbb{Z}^N}(1+|k|^2+|l|^2)^{-2N-1}\sum_{I\in\mathbb{D}_{\Gamma}}e_{k,l,I}(x)f_{k,l,I}(y).
\end{equation}
\end{enumerate}
\end{lemma}\begin{proof} We have
$$K(x,y)=\sum_{I\in\mathbb{D}_{\Gamma}}\chi_I(x)\chi_{3I}(y)K(x,y)\psi(\frac{\Phi(x)-\Phi(y)}{d_I})=$$
$$=\sum_{\substack{I\in\mathbb{D}_{\Gamma}\\ 3I\subset\Gamma}}\chi_I(x)\chi_{3I}(y)K(x,y)\psi(\frac{\Phi(x)-\Phi(y)}{d_I})+\sum_{\substack{I\in\mathbb{D}_{\Gamma}\\ 3I\not\subset\Gamma}}\chi_I(x)\chi_{3I}(y)K(x,y)\psi(\frac{\Phi(x)-\Phi(y)}{d_I}).$$

We will only consider the first sum. The second some can be treated similarly, but the notations become heavier. {\color{red} feel free to write the expression for the second sum if you think it is needed. the only difference is that the domain of the map $K_I$ becomes more complicated.}

For every $I\in\mathbb{D}_{\Gamma}$ with $3I\subset\Gamma,$ define the map
$$K_I:[-\frac12,\frac12]^N\times[-\frac32,\frac32]^N\to\mathbb{C}$$
by setting
$$K_I(a,b)=K(c_I+d_I\Phi^{-1}(a),c_I+d_I\Phi^{-1}(b))\psi(a-b),\quad a\in[-\frac12,\frac12]^N,\quad b\in[-\frac32,\frac32]^N.$$

It follows from the assumption on $K$ that {\color{red} feel free to prove this if you think this requires a proof. however, do not prove it here --- make a separate lemma instead.}
$$\|K_I\|_{C^{4N+2}([-\frac12,\frac12]^N\times[-\frac32,\frac32]^N)}\leq c_K\frac{d_I}{w(I)}.$$
Hence, there exists an extension $K_I:\mathbb{T}^N\times\mathbb{T}^N\to\mathbb{C}$ such that
$$\|K_I\|_{C^{4N+2}(\mathbb{T}^N\times\mathbb{T}^N)}\leq c_K'\frac{d_I}{w(I)}.$$
In particular,
$$\|K_I\|_{W^{4N+2,2}(\mathbb{T}^N\times\mathbb{T}^N)}\leq c_K''\frac{d_I}{w(I)}.$$
We may, therefore, write
$$K_I(a,b)=\sum_{k,l\in\mathbb{Z}^N}{\rm coef}_{k,l,I}e^{i\langle k,a\rangle}e^{i\langle l,b\rangle},\quad a\in[-\frac12,\frac12]^N,\quad b\in[-\frac32,\frac32]^N,$$
where
$$\sum_{k,l\in\mathbb{Z}^N}(1+|k|^2+|l|^2)^{4N+2}|{\rm coef}_{k,l,I}|^2\leq (c_K''\frac{d_I}{w(I)})^2.$$
In particular,
$$|{\rm coef}_{k,l,I}|\leq c_K''(1+|k|^2+|l|^2)^{-2N-1}\frac{d_I}{w(I)}.$$

This allows us to write
$$\sum_{\substack{I\in\mathbb{D}_{\Gamma}\\ 3I\subset\Gamma}}\chi_I(x)\chi_{3I}(y)K(x,y)\psi(\frac{\Phi(x)-\Phi(y)}{d_I})=\sum_{\substack{I\in\mathbb{D}_{\Gamma}\\ 3I\subset\Gamma}}\chi_I(x)\chi_{3I}(y)K_I(\frac{\Phi^{-1}(x-c_I)}{d_I},\frac{\Phi^{-1}(x)-c_I}{d_I})=$$
$$=\sum_{\substack{I\in\mathbb{D}_{\Gamma}\\ 3I\subset\Gamma}}\chi_I(x)\chi_{3I}(y)\sum_{k,l\in\mathbb{Z}^N}{\rm coef}_{k,l,I}e^{i\langle k,\frac{\Phi^{-1}(x-c_I)}{d_I}\rangle}e^{i\langle l,\frac{\Phi^{-1}(y-c_I)}{d_I}\rangle}=$$
$$=\sum_{k,l\in\mathbb{Z}^N}(1+|k|^2+|l|^2)^{-2N-1}\sum_{\substack{I\in\mathbb{D}_{\Gamma}\\ 3I\subset\Gamma}}(1+|k|^2+|l|^2)^{2N+1}{\rm coef}_{k,l,I}\cdot \chi_I(x)e^{i\langle k,\frac{\Phi^{-1}(x-c_I)}{d_I}\rangle}\cdot \chi_{3I}(y)e^{i\langle l,\frac{\Phi^{-1}(y-c_I)}{d_I}\rangle}.$$

We now set
$$e_{k,l,I}(x)=f(x)\chi_I(x)e^{i\langle k,\frac{\Phi^{-1}(x-c_I)}{d_I}\rangle},$$
$$f_{k,l,I}(y)=(1+|k|^2+|l|^2)^{2N+1}{\rm coef}_{k,l,I}\cdot \chi_{3I}(y)e^{i\langle l,\frac{\Phi^{-1}(y-c_I)}{d_I}\rangle}.$$

Obviously,
$${\rm supp}(e_{k,l,I})\subset I,\quad {\rm subset}(f_{k,l,I})\subset 3I,$$
$$\|e_{k,l,I}\|_{L_r(\mathbb{R}^N,w(x)dx)}=\|f\|_{L_r(I,w(x)dx)},$$
$$\|f_{k,l,I}\|_{L_r(\mathbb{R}^N,w(x)dx)}=(1+|k|^2+|l|^2)^{2N+1}|{\rm coef}_{I,k,l}|\cdot w(3I)^{\frac1r}\leq c_K''w(3I)^{\frac1r}\leq c_{K,R,k}w(I)^{\frac1r}.$$



\end{proof}

\begin{lemma}\label{lem:Cwikel-dyadic-distribution}
Let $1<s<\infty$, and let $\mathbb D^0$ be the standard dyadic grid
in $\mathbb R^N$. Suppose that
$
F\in L^s(\mathbb R^N),
$
For every $Q\in\mathbb D^0$, set
\[
b_Q
:=
|Q|^{\frac1s-1}\int_Q |F(x)|\,dx
=
|Q|^{-1/s'}\int_Q |F(x)|\,dx,
\]
Then
\begin{equation}\label{eq:Cwikel-dyadic-distribution}
\#\bigl\{Q\in\mathbb D^0:b_Q>t\bigr\}
\lesssim_{N,s}
t^{-s}\|F\|_{L^s(\mathbb R^N)}^s,
\qquad t>0.
\end{equation}
Equivalently,
\[
\left\|\{b_Q\}_{Q\in\mathbb D^0}\right\|_{\ell^{s,\infty}}
\lesssim_{N,s}
\|F\|_{L^s(\mathbb R^N)}.
\]
\end{lemma}

\begin{proof}
Let $\mathcal E\subset\mathbb D^0$ be an arbitrary finite collection
of dyadic cubes, and define
\[
G_{\mathcal E}(x)
:=
\sum_{Q\in\mathcal E}
|Q|^{-1/s'}\chi_Q(x).
\]
We first claim that
\begin{equation}\label{eq:Cwikel-dyadic-G}
\|G_{\mathcal E}\|_{L^{s'}(\mathbb R^N)}
\lesssim_{N,s}
(\#\mathcal E)^{1/s'}.
\end{equation}

Indeed, fix
\[
x\in\bigcup_{Q\in\mathcal E}Q.
\]
Set
\[
\mathcal E_x
:=
\{Q\in\mathcal E:x\in Q\}.
\]
Since $\mathcal E$ is finite, there exists a cube
$Q_x\in\mathcal E_x$ having minimal volume. Since any two cubes in
the same dyadic grid which contain the same point are nested, every
$Q\in\mathcal E_x$ contains $Q_x$. Thus the cubes in
$\mathcal E_x$ belong to the chain of dyadic ancestors of $Q_x$.

Denote by $Q_x^{(k)}$ the $k$-th dyadic ancestor of $Q_x$, so that
\[
|Q_x^{(k)}|
=
2^{kN}|Q_x|,
\qquad k\geq0.
\]
It follows that
\begin{align*}
G_{\mathcal E}(x)
&=
\sum_{Q\in\mathcal E_x}|Q|^{-1/s'}
\\
&\leq
\sum_{k=0}^{\infty}
|Q_x^{(k)}|^{-1/s'}
\\
&=
|Q_x|^{-1/s'}
\sum_{k=0}^{\infty}
2^{-kN/s'}
\\
&\lesssim_{N,s}
|Q_x|^{-1/s'}.
\end{align*}

For each $Q\in\mathcal E$, define
\[
E_Q
:=
\left\{
x\in\bigcup_{P\in\mathcal E}P:
Q_x=Q
\right\}.
\]
The sets $\{E_Q\}_{Q\in\mathcal E}$ are pairwise disjoint and
satisfy
\[
E_Q\subset Q.
\]
Therefore,
\begin{align*}
\|G_{\mathcal E}\|_{L^{s'}(\mathbb R^N)}^{s'}
&=
\int_{\mathbb R^N}
G_{\mathcal E}(x)^{s'}\,dx
\\
&\lesssim_{N,s}
\sum_{Q\in\mathcal E}
\int_{E_Q}|Q|^{-1}\,dx
\\
&=
\sum_{Q\in\mathcal E}
\frac{|E_Q|}{|Q|}
\\
&\leq
\sum_{Q\in\mathcal E}1
=
\#\mathcal E.
\end{align*}
Taking the $s'$-th root proves
\eqref{eq:Cwikel-dyadic-G}.

We now estimate the sum of $\{b_Q\}_{Q\in\mathcal E}$. By
H\"older's inequality and \eqref{eq:Cwikel-dyadic-G},
\begin{align*}
\sum_{Q\in\mathcal E}b_Q
&=
\sum_{Q\in\mathcal E}
|Q|^{-1/s'}
\int_Q|F(x)|\,dx
\\
&=
\int_{\mathbb R^N}
|F(x)|
\sum_{Q\in\mathcal E}
|Q|^{-1/s'}\chi_Q(x)\,dx
\\
&=
\int_{\mathbb R^N}
|F(x)|G_{\mathcal E}(x)\,dx
\\
&\leq
\|F\|_{L^s(\mathbb R^N)}
\|G_{\mathcal E}\|_{L^{s'}(\mathbb R^N)}
\\
&\lesssim_{N,s}
\|F\|_{L^s(\mathbb R^N)}
(\#\mathcal E)^{1/s'}.
\end{align*}

Now fix $t>0$ and let
\[
\mathcal E
\subset
\{Q\in\mathbb D^0:b_Q>t\}
\]
be an arbitrary finite subcollection. Since $b_Q>t$ for every
$Q\in\mathcal E$, we have
\[
t\,\#\mathcal E
<
\sum_{Q\in\mathcal E}b_Q
\lesssim_{N,s}
\|F\|_{L^s(\mathbb R^N)}
(\#\mathcal E)^{1/s'}.
\]
If $\mathcal E\neq\varnothing$, dividing by
$(\#\mathcal E)^{1/s'}$ and using
\[
1-\frac1{s'}=\frac1s,
\]
we obtain
\[
(\#\mathcal E)^{1/s}
\lesssim_{N,s}
t^{-1}\|F\|_{L^s(\mathbb R^N)}.
\]
Consequently,
\[
\#\mathcal E
\lesssim_{N,s}
t^{-s}\|F\|_{L^s(\mathbb R^N)}^s.
\]
Since this estimate holds for every finite subcollection
\[
\mathcal E
\subset
\{Q\in\mathbb D^0:b_Q>t\},
\]
it follows that
\[
\#\bigl\{Q\in\mathbb D^0:b_Q>t\bigr\}
\lesssim_{N,s}
t^{-s}\|F\|_{L^s(\mathbb R^N)}^s.
\]
This proves \eqref{eq:Cwikel-dyadic-distribution}, and hence also the
stated weak $\ell^s$ estimate.
\end{proof}

\begin{lemma}\label{lem:Cwikel-lambda-weak}
Assume that $N>2$ and $2<r<N$. Then, for every $\sigma\in G$ and
every $f\in L^N(\mathbb R^N,dx)$, we have
$f\in L^r_{\mathrm{loc}}(\mathbb R^N,dw)$ and
\begin{equation}\label{eq:Cwikel-lambda-weak}
\left\|
\left\{
\lambda_{\sigma,I}(f)
\right\}_{I\in\mathbb D_\Gamma}
\right\|_{\ell^{N,\infty}}
\lesssim
\|f\|_{L^N(\mathbb R^N,dx)}.
\end{equation}
The implicit constant is independent of $\sigma$ and $f$.
\end{lemma}
\begin{proof}
We first note that
$
f\in L^r_{\mathrm{loc}}(\mathbb R^N,dw).
$
Indeed, the density of $dw$ with respect to Lebesgue measure is
$
\prod_{v\in R_+}|\langle x,v\rangle|^{2k(v)},
$
which is bounded on every compact subset of $\mathbb R^N$. Hence, if
$E\subset\mathbb R^N$ is bounded, then, since $r<N$,
\[
\int_E |f(x)|^r\,dw(x)
\lesssim_E
\int_E |f(x)|^r\,dx
\leq
|E|^{1-\frac rN}
\|f\|_{L^N(\mathbb R^N,dx)}^r
<\infty.
\]
Thus $\lambda_{\sigma,I}(f)$ in
\eqref{eq:Dunkl-lambda-rank-one} is well defined.

Fix $\sigma\in G$. Recall that
$
\Omega
=
\mathbb R_+^{N_0}\times\mathbb R^{N-N_0}
$
and, by Lemma~\ref{mulemma},
\begin{equation}\label{eq:Cwikel-lambda-density}
d\mu(u,z)
=
\phi(u)\,du\,dz,
\qquad
\phi(u)
=
|{\rm det}(\Phi)|^{-1}
\prod_{v\in R_+}
\left(
\sum_{i=1}^{N_0}a_{v,i}u_i
\right)^{2k(v)},
\end{equation}
where
$
a_{v,i}\geq0.
$
In particular, $\phi$ is nondecreasing in each of the variables
$u_1,\ldots,u_{N_0}$ and is independent of the remaining
$N-N_0$ variables.

Let
$
I=\Phi^{-1}(Q)\in\mathbb D_\Gamma,
$
where
\[
Q
=
\prod_{i=1}^N
[m_i\ell,(m_i+1)\ell)
\subset\Omega,
\qquad
\ell=\ell(Q)=\ell(I),
\]
with
\[
m_i\in\mathbb Z_+,
\qquad 1\leq i\leq N_0,
\]
and
\[
m_i\in\mathbb Z,
\qquad N_0<i\leq N.
\]
Define
$
u^+
:=
\bigl(
(m_1+1)\ell,\ldots,(m_{N_0}+1)\ell
\bigr)
$
and
\[
Q^+
:=
\prod_{i=1}^{N_0}
\left[
\left(m_i+\frac12\right)\ell,
(m_i+1)\ell
\right)
\times
\prod_{i=N_0+1}^{N}
[m_i\ell,(m_i+1)\ell).
\]
Then
\[
|Q^+|=2^{-N_0}|Q|.
\]

For every $(u,z)\in Q^+$ and every $1\leq i\leq N_0$,
\[
u_i
\geq
\left(m_i+\frac12\right)\ell
\geq
\frac12(m_i+1)\ell.
\]
Since $a_{v,i}\geq0$, it follows that, for every $v\in R_+$,
\[
\sum_{i=1}^{N_0}a_{v,i}u_i
\geq
\frac12
\sum_{i=1}^{N_0}a_{v,i}(m_i+1)\ell.
\]
Consequently,
\[
\phi(u)
\geq
2^{-2\sum_{v\in R_+}k(v)}
\phi(u^+).
\]
Since
\[
2\sum_{v\in R_+}k(v)
=
\sum_{v\in R}k(v)
=
\N-N,
\]
we obtain
\[
\phi(u)
\geq
2^{-(\N-N)}
\phi(u^+),
\qquad
(u,z)\in Q^+.
\]
Therefore,
\begin{align*}
\mu(Q)
&=
\int_Q\phi(u)\,du\,dz
\\
&\geq
\int_{Q^+}\phi(u)\,du\,dz
\\
&\geq
2^{-(\N-N)}
|Q^+|\phi(u^+)
\\
&=
2^{-(\N-N+N_0)}
|Q|\phi(u^+).
\end{align*}
On the other hand, by the coordinatewise monotonicity of $\phi$,
\[
\operatorname*{ess\,sup}_{(u,z)\in Q}\phi(u)
\leq
\phi(u^+).
\]
It follows that
\begin{equation}\label{eq:Cwikel-density-sup-average}
\operatorname*{ess\,sup}_{Q}\phi
\lesssim
\frac{\mu(Q)}{|Q|}.
\end{equation}

Since $dw$ is $G$-invariant and
\[
\mu(Q)=w(\Phi^{-1}(Q))=w(I),
\]
we have
\begin{align*}
\frac{1}{w(I)}
\int_{\sigma I}|f(x)|^r\,dw(x)
&=
\frac{1}{w(I)}
\int_I|f(\sigma x)|^r\,dw(x)
\\
&=
\frac{1}{\mu(Q)}
\int_Q
|f(\sigma\Phi^{-1}(\xi))|^r\,d\mu(\xi)
\\
&=
\frac{1}{\mu(Q)}
\int_Q
|f(\sigma\Phi^{-1}(\xi))|^r
\phi(\xi)\,d\xi
\\
&\leq
\frac{\operatorname*{ess\,sup}_{Q}\phi}{\mu(Q)}
\int_Q
|f(\sigma\Phi^{-1}(\xi))|^r\,d\xi
\\
&\lesssim
\frac1{|Q|}
\int_Q
|f(\sigma\Phi^{-1}(\xi))|^r\,d\xi,
\end{align*}
where we used \eqref{eq:Cwikel-density-sup-average} in the last
inequality. Hence, by the definition
\eqref{eq:Dunkl-lambda-rank-one},
\begin{align}
\lambda_{\sigma,I}(f)^r
&=
\ell(I)^r
\frac1{w(I)}
\int_{\sigma I}|f(x)|^r\,dw(x)
\nonumber\\
&\lesssim
\frac{\ell(Q)^r}{|Q|}
\int_Q
|f(\sigma\Phi^{-1}(\xi))|^r\,d\xi.
\label{eq:Cwikel-lambda-Lebesgue}
\end{align}

Set
\[
s:=\frac Nr.
\]
Since $r<N$, we have $s>1$. Moreover,
\[
|Q|=\ell(Q)^N,
\]
and hence
\[
\frac{\ell(Q)^r}{|Q|}
=
|Q|^{\frac rN-1}
=
|Q|^{\frac1s-1}.
\]
Define
\[
F_\sigma(\xi)
:=
\begin{cases}
|f(\sigma\Phi^{-1}(\xi))|^r,
&\xi\in\Omega,\\
0,
&\xi\notin\Omega.
\end{cases}
\]
Then \eqref{eq:Cwikel-lambda-Lebesgue} gives
\begin{equation}\label{eq:Cwikel-lambda-dyadic}
\lambda_{\sigma,I}(f)^r
\lesssim
|Q|^{\frac1s-1}
\int_QF_\sigma(\xi)\,d\xi.
\end{equation}

We next verify that $F_\sigma\in L^s(\mathbb R^N)$. Since
$
rs=N,
$
a change of variables $\xi=\Phi(x)$ gives
\begin{align*}
\|F_\sigma\|_{L^s(\mathbb R^N)}^s
&=
\int_\Omega
|f(\sigma\Phi^{-1}(\xi))|^{rs}\,d\xi
\\
&=
\int_\Omega
|f(\sigma\Phi^{-1}(\xi))|^N\,d\xi
\\
&=
|{\rm det}(\Phi)|
\int_\Gamma
|f(\sigma x)|^N\,dx
\\
&=
|{\rm det}(\Phi)|
\int_{\sigma\Gamma}
|f(x)|^N\,dx
\\
&\leq
|{\rm det}(\Phi)|
\|f\|_{L^N(\mathbb R^N,dx)}^N
<\infty,
\end{align*}
where in the fourth line we used the fact that $\sigma$ is
orthogonal. Thus
$
F_\sigma\in L^s(\mathbb R^N).
$

Since $\mathbb D_\Omega^0$ is a subfamily of the standard dyadic
system $\mathbb D^0$, we may apply
Lemma~\ref{lem:Cwikel-dyadic-distribution} to $F_\sigma$. Together
with \eqref{eq:Cwikel-lambda-dyadic}, this gives, for every $a>0$,
\begin{align*}
\#\left\{
I\in\mathbb D_\Gamma:
\lambda_{\sigma,I}(f)>a
\right\}
&\leq
\#\left\{
Q\in\mathbb D_\Omega^0:
|Q|^{\frac1s-1}
\int_QF_\sigma(\xi)\,d\xi
\gtrsim a^r
\right\}
\\
&\leq
\#\left\{
Q\in\mathbb D^0:
|Q|^{\frac1s-1}
\int_QF_\sigma(\xi)\,d\xi
\gtrsim a^r
\right\}
\\
&\lesssim
(a^r)^{-s}
\|F_\sigma\|_{L^s(\mathbb R^N)}^s
\\
&=
a^{-rs}
\|F_\sigma\|_{L^s(\mathbb R^N)}^s
\\
&\lesssim
a^{-N}
\|f\|_{L^N(\mathbb R^N,dx)}^N,
\end{align*}
where in the last line we used $rs=N$ and the estimate for
$\|F_\sigma\|_{L^s(\mathbb R^N)}^s$ obtained above. Therefore,
\[
\#\left\{
I\in\mathbb D_\Gamma:
\lambda_{\sigma,I}(f)>a
\right\}
\lesssim
\left(
\frac{\|f\|_{L^N(\mathbb R^N,dx)}}{a}
\right)^N,
\qquad a>0.
\]
By the definition of the weak $\ell^N$ quasi-norm, this is exactly
\[
\left\|
\left\{
\lambda_{\sigma,I}(f)
\right\}_{I\in\mathbb D_\Gamma}
\right\|_{\ell^{N,\infty}}
\lesssim
\|f\|_{L^N(\mathbb R^N,dx)}.
\]
All the implicit constants above depend only on the fixed structural
parameters and are independent of $\sigma\in G$ and $f$. This
completes the proof.
\end{proof}
\begin{lemma}\label{lem:Cwikel-rank-one-block}
Assume that $N>2$ and $2<r<N$.
Let
\[
\left\{
e_{p,q,I}^{\sigma,\tau}
\right\}_{p,q\in\mathbb Z^N,\,
I\in\mathbb D_\Gamma,\,
\sigma,\tau\in G}
\quad\text{and}\quad
\left\{
h_{p,q,I}^{\sigma,\tau}
\right\}_{p,q\in\mathbb Z^N,\,
I\in\mathbb D_\Gamma,\,
\sigma,\tau\in G}
\]
be the families obtained from
Lemma~\ref{lem:rank-one-decomposition-Dunkl} with this choice of $r$.
For every
\[
f\in L^N(\mathbb R^N,dx),
\qquad
\sigma,\tau\in G,
\qquad
p,q\in\mathbb Z^N,
\]
define
\[
A_{\sigma,\tau,p,q}(f)
:=
\sum_{I\in\mathbb D_\Gamma}
\lambda_{\sigma,I}(f)
\left\langle
\,\cdot,\,
\overline{h_{p,q,I}^{\sigma,\tau}}
\right\rangle_{L^2(\mathbb R^N,dw)}
e_{p,q,I}^{\sigma,\tau}.
\]
Then
\[
A_{\sigma,\tau,p,q}(f)
\in
S^{N,\infty}\bigl(L^2(\mathbb R^N,dw)\bigr)
\]
and
\begin{equation}\label{eq:Cwikel-rank-one-block}
\left\|
A_{\sigma,\tau,p,q}(f)
\right\|_
{S^{N,\infty}(L^2(\mathbb R^N,dw))}
\lesssim
\|f\|_{L^N(\mathbb R^N,dx)},
\end{equation}
where the implicit constant is independent of
$\sigma,\tau,p,q$, and $f$.
\end{lemma}

\begin{proof}
By Lemma~\ref{lem:Cwikel-lambda-weak},
\begin{equation}\label{eq:Cwikel-block-lambda-weak}
\left\|
\left\{
\lambda_{\sigma,I}(f)
\right\}_{I\in\mathbb D_\Gamma}
\right\|_{\ell^{N,\infty}}
\lesssim
\|f\|_{L^N(\mathbb R^N,dx)},
\end{equation}
uniformly in $\sigma\in G$.

Fix
\[
\sigma,\tau\in G,
\qquad
p,q\in\mathbb Z^N.
\]
For $\gamma\in G$, define
\[
U_\gamma:
L^2(\Gamma,dw)
\longrightarrow
L^2(\gamma\Gamma,dw)
\]
by
\[
U_\gamma u(x)
:=
u(\gamma^{-1}x),
\qquad x\in\gamma\Gamma.
\]
Since $dw$ is $G$-invariant, $U_\gamma$ is unitary. In fact, the same
argument shows that $U_\gamma$ is an isometry on $L^s(dw)$ for every
$1\leq s\leq\infty$.

For $I\in\mathbb D_\Gamma$, define
\[
\widetilde e_I
:=
U_\sigma^{-1}
e_{p,q,I}^{\sigma,\tau}
\]
and
\[
\widetilde h_I
:=
U_\tau^{-1}
\overline{h_{p,q,I}^{\sigma,\tau}}.
\]
By \eqref{eq:Dunkl-support-e-new},
\[
\operatorname{supp}\widetilde e_I
\subset I,
\qquad
\operatorname{supp}\widetilde h_I
\subset I^\ast.
\]

We claim that
\[
\{\widetilde e_I\}_{I\in\mathbb D_\Gamma}
\quad\text{and}\quad
\{\widetilde h_I\}_{I\in\mathbb D_\Gamma}
\]
are NWO sequences.

Indeed, write
\[
I=\Phi^{-1}(Q),
\qquad
I^\ast=\Phi^{-1}\bigl((2Q)\cap\Omega\bigr).
\]
Since $\Phi$ and $\Phi^{-1}$ are bi-Lipschitz, the set $I^\ast$ is
contained in a fixed geometric enlargement of $I$, with the enlargement
constant independent of $I$. Thus the support condition in the
definition of an NWO sequence is satisfied for both families.

Moreover, since $U_\sigma$ and $U_\tau$ preserve the $L^r(dw)$ norm,
\eqref{eq:Dunkl-normalization-e-new} gives
\begin{equation}\label{eq:Cwikel-block-NWO-norm-e}
\|\widetilde e_I\|_{L^r(\Gamma,dw)}
=
\|e_{p,q,I}^{\sigma,\tau}\|_{L^r(\mathbb R^N,dw)}
\leq
w(I)^{\frac1r-\frac12}
\end{equation}
and
\begin{equation}\label{eq:Cwikel-block-NWO-norm-h}
\|\widetilde h_I\|_{L^r(\Gamma,dw)}
=
\|h_{p,q,I}^{\sigma,\tau}\|_{L^r(\mathbb R^N,dw)}
\leq
w(I)^{\frac1r-\frac12}.
\end{equation}
Since $r>2$, these are precisely the required $L^r$ normalization
conditions for NWO sequences. Therefore
\[
\{\widetilde e_I\}_{I\in\mathbb D_\Gamma}
\quad\text{and}\quad
\{\widetilde h_I\}_{I\in\mathbb D_\Gamma}
\]
are NWO sequences.

Consider now the operator on $L^2(\Gamma,dw)$ given by
\begin{equation}\label{eq:Cwikel-pullback-block}
\widetilde A_{\sigma,\tau,p,q}(f)
:=
\sum_{I\in\mathbb D_\Gamma}
\lambda_{\sigma,I}(f)
\left\langle
\,\cdot,\,
\widetilde h_I
\right\rangle_{L^2(\Gamma,dw)}
\widetilde e_I.
\end{equation}
Applying Lemma~\ref{NWOLEMMAA2} with exponent $N$, and then
using \eqref{eq:Cwikel-block-lambda-weak}, we obtain
\begin{align}
\left\|
\widetilde A_{\sigma,\tau,p,q}(f)
\right\|_
{S^{N,\infty}(L^2(\Gamma,dw))}
&\lesssim
\left\|
\left\{
\lambda_{\sigma,I}(f)
\right\}_{I\in\mathbb D_\Gamma}
\right\|_{\ell^{N,\infty}}\lesssim
\|f\|_{L^N(\mathbb R^N,dx)}.
\label{eq:Cwikel-pullback-block-estimate}
\end{align}

It remains to transfer this estimate back to
$L^2(\mathbb R^N,dw)$.

By \eqref{eq:Dunkl-support-e-new}, the range of
$A_{\sigma,\tau,p,q}(f)$ is contained in
$L^2(\sigma\Gamma,dw)$, whereas the operator vanishes on
$L^2(\mathbb R^N\setminus\tau\Gamma,dw)$. Thus its nonzero part is an
operator
\[
L^2(\tau\Gamma,dw)
\longrightarrow
L^2(\sigma\Gamma,dw).
\]
For $g\in L^2(\tau\Gamma,dw)$, put
\[
u:=U_\tau^{-1}g.
\]
Then, using the definition of $\widetilde h_I$ and the unitarity of
$U_\tau$,
\begin{align*}
\left\langle
g,
\overline{h_{p,q,I}^{\sigma,\tau}}
\right\rangle_{L^2(\tau\Gamma,dw)}
&=
\left\langle
U_\tau^{-1}g,
U_\tau^{-1}
\overline{h_{p,q,I}^{\sigma,\tau}}
\right\rangle_{L^2(\Gamma,dw)}
=
\langle u,\widetilde h_I\rangle_{L^2(\Gamma,dw)}.
\end{align*}
Also,
\[
e_{p,q,I}^{\sigma,\tau}
=
U_\sigma\widetilde e_I.
\]
Consequently,
\begin{align*}
A_{\sigma,\tau,p,q}(f)g
&=
\sum_{I\in\mathbb D_\Gamma}
\lambda_{\sigma,I}(f)
\left\langle
g,
\overline{h_{p,q,I}^{\sigma,\tau}}
\right\rangle
e_{p,q,I}^{\sigma,\tau}
\\
&=
U_\sigma
\sum_{I\in\mathbb D_\Gamma}
\lambda_{\sigma,I}(f)
\langle u,\widetilde h_I\rangle
\widetilde e_I
\\
&=
U_\sigma
\widetilde A_{\sigma,\tau,p,q}(f)
U_\tau^{-1}g.
\end{align*}
Hence
\begin{equation}\label{eq:Cwikel-block-unitary-equivalence}
A_{\sigma,\tau,p,q}(f)
=
U_\sigma
\widetilde A_{\sigma,\tau,p,q}(f)
U_\tau^{-1}
\end{equation}
on $L^2(\tau\Gamma,dw)$, with the left-hand side extended by zero on
the orthogonal complement of $L^2(\tau\Gamma,dw)$.

Therefore the nonzero singular values of
$A_{\sigma,\tau,p,q}(f)$ and
$\widetilde A_{\sigma,\tau,p,q}(f)$ coincide. In particular,
\[
\left\|
A_{\sigma,\tau,p,q}(f)
\right\|_{S^{N,\infty}(L^2(\mathbb R^N,dw))}
=
\left\|
\widetilde A_{\sigma,\tau,p,q}(f)
\right\|_{S^{N,\infty}(L^2(\Gamma,dw))}.
\]
Combining this identity with
\eqref{eq:Cwikel-pullback-block-estimate}, we conclude that
\[
\left\|
A_{\sigma,\tau,p,q}(f)
\right\|_
{S^{N,\infty}(L^2(\mathbb R^N,dw))}
\lesssim
\|f\|_{L^N(\mathbb R^N,dx)}.
\]
The implicit constant is independent of
$\sigma,\tau,p,q$, and $f$, which proves
\eqref{eq:Cwikel-rank-one-block}.
\end{proof}

\begin{lemma}\label{lem:Cwikel-Dunkl}
Assume that
$
N>2$.
Then there exists a constant $C_{N,R,k}>0$ such that, for every
$f\in L^N(\mathbb R^N,dx)$,
\[
M_f(-\Delta_{\D})^{-1/2}
\in
S^{N,\infty}\bigl(L^2(\mathbb R^N,dw)\bigr)
\]
and
\begin{equation}\label{eq:Cwikel-Dunkl}
\bigl\|
M_f(-\Delta_{\D})^{-1/2}
\bigr\|_
{S^{N,\infty}(L^2(\mathbb R^N,dw))}
\leq
C_{N,R,k}
\|f\|_{L^N(\mathbb R^N,dx)}.
\end{equation}
\end{lemma}
\begin{proof}
Choose
$
2<r<N.
$
By Lemma~\ref{lem:Cwikel-lambda-weak},
$
f\in L^r_{\mathrm{loc}}(\mathbb R^N,dw),
$
and hence Lemma~\ref{lem:rank-one-decomposition-Dunkl} is applicable
with this choice of $r$.

Let
\[
\{c_{p,q}\}_{p,q\in\mathbb Z^N},
\qquad
\{e_{p,q,I}^{\sigma,\tau}\},
\qquad
\{h_{p,q,I}^{\sigma,\tau}\}
\]
be the coefficients and functions given by
Lemma~\ref{lem:rank-one-decomposition-Dunkl}. For
$\sigma,\tau\in G$ and $p,q\in\mathbb Z^N$, let
\[
A_{\sigma,\tau,p,q}(f)
:=
\sum_{I\in\mathbb D_\Gamma}
\lambda_{\sigma,I}(f)
\left\langle
\,\cdot,\,
\overline{h_{p,q,I}^{\sigma,\tau}}
\right\rangle_{L^2(\mathbb R^N,dw)}
e_{p,q,I}^{\sigma,\tau}.
\]
By Lemma~\ref{lem:Cwikel-rank-one-block},
\begin{equation}\label{eq:Cwikel-block-estimate-final}
\left\|
A_{\sigma,\tau,p,q}(f)
\right\|_
{S^{N,\infty}(L^2(\mathbb R^N,dw))}
\lesssim
\|f\|_{L^N(\mathbb R^N,dx)}
\end{equation}
uniformly in
\[
\sigma,\tau\in G,
\qquad
p,q\in\mathbb Z^N.
\]

On the other hand, by \eqref{eq:Dunkl-cpq-decay}, for every $M>0$,
\[
c_{p,q}
\leq
C_M(1+|p|+|q|)^{-M}.
\]
Taking $M>2N$, we obtain
\begin{equation}\label{eq:Cwikel-cpq-l1-final}
\sum_{p,q\in\mathbb Z^N}c_{p,q}<\infty.
\end{equation}
Since $G$ is finite, \eqref{eq:Cwikel-block-estimate-final} and
\eqref{eq:Cwikel-cpq-l1-final} imply
\begin{align*}
&\sum_{\sigma,\tau\in G}
\sum_{p,q\in\mathbb Z^N}
c_{p,q}
\left\|
A_{\sigma,\tau,p,q}(f)
\right\|_
{S^{N,\infty}(L^2(\mathbb R^N,dw))}
\\
&\qquad\lesssim
|G|^2
\left(
\sum_{p,q\in\mathbb Z^N}c_{p,q}
\right)
\|f\|_{L^N(\mathbb R^N,dx)}
\\
&\qquad\lesssim_{N,R,k}
\|f\|_{L^N(\mathbb R^N,dx)}.
\end{align*}
Therefore the series
\[
\sum_{\sigma,\tau\in G}
\sum_{p,q\in\mathbb Z^N}
c_{p,q}A_{\sigma,\tau,p,q}(f)
\]
converges in
$S^{N,\infty}(L^2(\mathbb R^N,dw))$, and
\begin{equation}\label{eq:Cwikel-summed-blocks}
\left\|
\sum_{\sigma,\tau\in G}
\sum_{p,q\in\mathbb Z^N}
c_{p,q}A_{\sigma,\tau,p,q}(f)
\right\|_
{S^{N,\infty}(L^2(\mathbb R^N,dw))}
\lesssim_{N,R,k}
\|f\|_{L^N(\mathbb R^N,dx)}.
\end{equation}

It remains only to identify the operator represented by this series.
By the definition of $A_{\sigma,\tau,p,q}(f)$, the integral kernel of
its summand corresponding to $I\in\mathbb D_\Gamma$ is
\[
\lambda_{\sigma,I}(f)
e_{p,q,I}^{\sigma,\tau}(x)
h_{p,q,I}^{\sigma,\tau}(y).
\]
Hence the kernel identity
\eqref{eq:Dunkl-kernel-rank-one-new} in
Lemma~\ref{lem:rank-one-decomposition-Dunkl} gives
\[
M_f(-\Delta_{\D})^{-1/2}
=
C
\sum_{\sigma,\tau\in G}
\sum_{p,q\in\mathbb Z^N}
c_{p,q}A_{\sigma,\tau,p,q}(f).
\]
More precisely, the two sides have the same integral kernel
\[
f(x)K_{(-\Delta_{\D})^{-1/2}}(x,y)
\]
for almost every $(x,y)$ with $d(x,y)>0$; hence they agree initially
on the natural test-function class and therefore, by the boundedness
of the right-hand side, as operators on $L^2(\mathbb R^N,dw)$.

Combining this identity with \eqref{eq:Cwikel-summed-blocks}, we obtain
\[
M_f(-\Delta_{\D})^{-1/2}
\in
S^{N,\infty}\bigl(L^2(\mathbb R^N,dw)\bigr)
\]
and
\[
\left\|
M_f(-\Delta_{\D})^{-1/2}
\right\|_
{S^{N,\infty}(L^2(\mathbb R^N,dw))}
\lesssim_{N,R,k}
\|f\|_{L^N(\mathbb R^N,dx)}.
\]
Absorbing the structural constants into $C_{N,R,k}$ gives
\eqref{eq:Cwikel-Dunkl}.
\end{proof}

%
%\begin{proof}
%Fix once and for all an exponent
%\[
%s>2.
%\]
%For instance, one may take \(s=4\). Since
%\(f\in\dot W^{1,N}(\mathbb R^N)\), the local critical
%Sobolev--Poincar\'e inequality implies
%\[
%f\in L^s_{\rm loc}(\mathbb R^N,dx).
%\]
%Since the density \(dw/dx\) is bounded on compact subsets,
%\[
%f\in L^s_{\rm loc}(\mathbb R^N,dw).
%\]
%Thus all the local expressions below are well defined.
%
%We use the endpoint decomposition from
%Lemma~\ref{lem:Dunkl-Riesz-Fourier-NWO-endpoint}. Suppressing
%temporarily the indices \(u,v,\sigma,\tau\), write
%\[
%e_I=e_{u,v,I}^{\sigma,\tau},
%\qquad
%h_I=h_{u,v,I}^{\sigma,\tau},
%\qquad
%\theta_I=\theta_{\sigma,\tau,I}.
%\]
%On the support of \(e_I(x)h_I(y)\), we have
%\[
%x\in\sigma I,
%\qquad
%y\in\tau I^*.
%\]
%Set
%\[
%\xi:=\sigma^{-1}x\in I,
%\qquad
%\eta:=\tau^{-1}y\in I^*.
%\]
%Then
%\begin{align}
%f(y)-f(x)
%&=
%f(\tau\eta)-f(\sigma\xi)
%\nonumber\\
%&=
%\bigl[
%f(\tau\eta)-f(\tau\xi)
%\bigr]
%+
%\bigl[
%f(\tau\xi)-f(\sigma\xi)
%\bigr].
%\label{eq:commutator-two-parts}
%\end{align}
%We treat these two terms separately.
%
%\medskip
%\noindent
%\textbf{Step 1. The same-chamber term.}
%
%Let
%\[
%g_\tau(z):=f(\tau z),
%\qquad z\in\Gamma.
%\]
%Let \(\mathcal B_I\) be the chamber ball from
%Lemma~\ref{lem:Dunkl-Ainfty-critical-oscillation}, and put
%\[
%m_{\tau,I}
%:=
%(g_\tau)_{\mathcal B_I,dx}.
%\]
%Define
%\begin{equation}
%\label{eq:local-same-chamber-omega}
%\omega_{\tau,I}
%:=
%\left[
%\frac1{w(I)}
%\int_{I^*}
%|g_\tau(z)-m_{\tau,I}|^s\,dw(z)
%\right]^{1/s}.
%\end{equation}
%Since
%\[
%I^*\subset\mathcal B_I,
%\qquad
%w(\mathcal B_I)\simeq w(I),
%\]
%Lemma~\ref{lem:Dunkl-Ainfty-critical-oscillation} gives
%\begin{equation}
%\label{eq:local-omega-weak}
%\left\|
%\{\omega_{\tau,I}\}_{I\in\mathbb D_\Gamma}
%\right\|_{\ell^{N,\infty}}
%\lesssim
%\|\nabla g_\tau\|_{L^N(\Gamma)}
%\leq
%\|\nabla f\|_{L^N(\mathbb R^N)}.
%\end{equation}
%Since \(0<\theta_I\leq1\),
%\begin{equation}
%\label{eq:theta-local-omega-weak}
%\left\|
%\{\theta_I\omega_{\tau,I}\}_{I}
%\right\|_{\ell^{N,\infty}}
%\lesssim
%\|\nabla f\|_{L^N}.
%\end{equation}
%
%If \(\omega_{\tau,I}>0\), define
%\[
%E_{I,0}:=e_I,
%\qquad
%H_{I,0}(y)
%:=
%\omega_{\tau,I}^{-1}
%\bigl(f(y)-m_{\tau,I}\bigr)h_I(y),
%\]
%and
%\[
%E_{I,1}(x)
%:=
%-\omega_{\tau,I}^{-1}
%\bigl(
%f(\tau\sigma^{-1}x)-m_{\tau,I}
%\bigr)e_I(x),
%\qquad
%H_{I,1}:=h_I.
%\]
%If \(\omega_{\tau,I}=0\), set these modified factors equal to zero.
%Then
%\begin{align}
%&
%\theta_I
%\bigl[
%f(\tau\eta)-f(\tau\xi)
%\bigr]
%e_I(x)h_I(y)
%\nonumber\\
%&\qquad=
%\theta_I\omega_{\tau,I}
%\sum_{\varepsilon=0}^{1}
%E_{I,\varepsilon}(x)H_{I,\varepsilon}(y).
%\label{eq:same-chamber-NWO-split}
%\end{align}
%
%We claim that, for each \(\varepsilon=0,1\),
%\[
%\{E_{I,\varepsilon}\}_{I\in\mathbb D_\Gamma},
%\qquad
%\{H_{I,\varepsilon}\}_{I\in\mathbb D_\Gamma}
%\]
%are NWO families after the same \(G\)-pullbacks as in
%Lemma~\ref{lem:Dunkl-Riesz-Fourier-NWO-endpoint}.
%The support conditions are immediate, since multiplication by the
%displayed functions does not enlarge supports.
%
%For example, using
%\eqref{eq:endpoint-e-h-pointwise},
%\begin{align*}
%\|H_{I,0}\|_{L^s(dw)}
%&\lesssim
%\omega_{\tau,I}^{-1}
%w(I)^{-1/2}
%\left(
%\int_{\tau I^*}
%|f-m_{\tau,I}|^s\,dw
%\right)^{1/s}
%\\
%&=
%\omega_{\tau,I}^{-1}
%w(I)^{-1/2}
%\left(
%\int_{I^*}
%|g_\tau-m_{\tau,I}|^s\,dw
%\right)^{1/s}
%\\
%&=
%w(I)^{1/s-1/2}.
%\end{align*}
%Similarly, by the \(G\)-invariance of \(dw\),
%\begin{align*}
%\|E_{I,1}\|_{L^s(dw)}
%&\lesssim
%\omega_{\tau,I}^{-1}
%w(I)^{-1/2}
%\left(
%\int_{\sigma I}
%|f(\tau\sigma^{-1}x)-m_{\tau,I}|^s\,dw(x)
%\right)^{1/s}
%\\
%&=
%\omega_{\tau,I}^{-1}
%w(I)^{-1/2}
%\left(
%\int_I
%|g_\tau(z)-m_{\tau,I}|^s\,dw(z)
%\right)^{1/s}
%\\
%&\leq
%w(I)^{1/s-1/2}.
%\end{align*}
%The other two factors are the original NWO factors. Hence the claim
%follows.
%
%\medskip
%\noindent
%\textbf{Step 2. The cross-chamber term.}
%
%Define
%\[
%H_{\sigma,\tau}(z)
%=
%f(\tau z)-f(\sigma z),
%\qquad z\in\Gamma,
%\]
%and
%\[
%a_{\sigma,\tau,I}
%=
%\frac1{|I|}
%\int_I H_{\sigma,\tau}(z)\,dz.
%\]
%Put
%\begin{equation}
%\label{eq:jump-residual-rho}
%\rho_{\sigma,\tau,I}
%:=
%\left[
%\frac1{w(I)}
%\int_I
%|H_{\sigma,\tau}(z)-a_{\sigma,\tau,I}|^s\,dw(z)
%\right]^{1/s}.
%\end{equation}
%By Lemma~\ref{lem:Dunkl-Ainfty-critical-oscillation},
%\begin{align}
%\left\|
%\{\rho_{\sigma,\tau,I}\}_{I}
%\right\|_{\ell^{N,\infty}}
%&\lesssim
%\|\nabla H_{\sigma,\tau}\|_{L^N(\Gamma)}
%\nonumber\\
%&\lesssim
%\|\nabla f\|_{L^N(\mathbb R^N)}.
%\label{eq:jump-rho-weak}
%\end{align}
%Indeed,
%\[
%\nabla H_{\sigma,\tau}(z)
%=
%\tau^T\nabla f(\tau z)
%-
%\sigma^T\nabla f(\sigma z),
%\]
%and therefore
%\[
%\|\nabla H_{\sigma,\tau}\|_{L^N(\Gamma)}
%\leq
%2\|\nabla f\|_{L^N(\mathbb R^N)}.
%\]
%Consequently,
%\begin{equation}
%\label{eq:theta-jump-rho-weak}
%\left\|
%\{\theta_I\rho_{\sigma,\tau,I}\}_{I}
%\right\|_{\ell^{N,\infty}}
%\lesssim
%\|\nabla f\|_{L^N}.
%\end{equation}
%
%If \(\rho_{\sigma,\tau,I}>0\), define
%\[
%E_{I,2}(x)
%:=
%\rho_{\sigma,\tau,I}^{-1}
%\bigl(
%H_{\sigma,\tau}(\sigma^{-1}x)
%-
%a_{\sigma,\tau,I}
%\bigr)e_I(x),
%\qquad
%H_{I,2}:=h_I.
%\]
%If \(\rho_{\sigma,\tau,I}=0\), put \(E_{I,2}=0\).
%Then
%\begin{align}
%&
%\theta_I
%H_{\sigma,\tau}(\xi)e_I(x)h_I(y)
%\nonumber\\
%&\qquad=
%\theta_I\rho_{\sigma,\tau,I}
%E_{I,2}(x)H_{I,2}(y)
%+
%\theta_Ia_{\sigma,\tau,I}
%e_I(x)h_I(y).
%\label{eq:cross-chamber-NWO-split}
%\end{align}
%Exactly as above,
%\[
%\|E_{I,2}\|_{L^s(dw)}
%\lesssim
%w(I)^{1/s-1/2},
%\]
%because
%\begin{align*}
%&
%\int_{\sigma I}
%\left|
%H_{\sigma,\tau}(\sigma^{-1}x)
%-a_{\sigma,\tau,I}
%\right|^s\,dw(x)
%\\
%&\qquad=
%\int_I
%|H_{\sigma,\tau}(z)-a_{\sigma,\tau,I}|^s\,dw(z).
%\end{align*}
%Thus
%\(\{E_{I,2}\}_I\) and \(\{H_{I,2}\}_I\) are NWO families.
%
%The coefficient of the remaining scalar term in
%\eqref{eq:cross-chamber-NWO-split} is controlled by
%Lemma~\ref{lem:Dunkl-chamber-jump-averages}:
%\begin{equation}
%\label{eq:scalar-jump-weak-final}
%\left\|
%\left\{
%\theta_Ia_{\sigma,\tau,I}
%\right\}_{I}
%\right\|_{\ell^{N,\infty}}
%\lesssim
%\|\nabla f\|_{L^N}.
%\end{equation}
%
%\medskip
%\noindent
%\textbf{Step 3. Application of the NWO--Schatten estimate.}
%
%Combining
%\eqref{eq:commutator-two-parts},
%\eqref{eq:same-chamber-NWO-split}, and
%\eqref{eq:cross-chamber-NWO-split}, we obtain, on each rank-one
%piece of
%\eqref{eq:endpoint-Dunkl-Riesz-kernel-NWO},
%\begin{align}
%&
%\theta_I
%(f(y)-f(x))
%e_I(x)h_I(y)
%\nonumber\\
%&\quad=
%\theta_I\omega_{\tau,I}
%\sum_{\varepsilon=0}^{1}
%E_{I,\varepsilon}(x)H_{I,\varepsilon}(y)
%\nonumber\\
%&\qquad+
%\theta_I\rho_{\sigma,\tau,I}
%E_{I,2}(x)H_{I,2}(y)
%+
%\theta_Ia_{\sigma,\tau,I}
%e_I(x)h_I(y).
%\label{eq:four-NWO-pieces}
%\end{align}
%
%For each fixed
%\[
%\sigma,\tau\in G,
%\qquad
%u,v\in\mathbb Z^N,
%\]
%all the families appearing on the right-hand side of
%\eqref{eq:four-NWO-pieces} are NWO families. Therefore
%Lemma~\ref{NWOLEMMAA2}, with
%\[
%p=N,
%\qquad
%q=\infty,
%\]
%together with
%\eqref{eq:theta-local-omega-weak},
%\eqref{eq:theta-jump-rho-weak}, and
%\eqref{eq:scalar-jump-weak-final}, gives
%\begin{align}
%&
%\left\|
%\sum_{I\in\mathbb D_\Gamma}
%\theta_I(f(y)-f(x))
%e_I(x)h_I(y)
%\right\|_{S^{N,\infty}}
%\nonumber\\
%&\qquad\lesssim
%\|\nabla f\|_{L^N(\mathbb R^N)},
%\label{eq:one-frequency-Schatten}
%\end{align}
%where, as usual, the displayed kernel is understood as the
%corresponding sum of rank-one operators
%\[
%\lambda_I
%\langle\,\cdot,\overline{H_I}\rangle F_I.
%\]
%The constant in
%\eqref{eq:one-frequency-Schatten} is uniform in
%\(\sigma,\tau,u,v\).
%
%We now return to
%\eqref{eq:endpoint-Dunkl-Riesz-kernel-NWO}. Since
%\(M_0>2N\),
%\[
%\sum_{u,v\in\mathbb Z^N}|c_{u,v}|<\infty.
%\]
%Moreover, \(G\) is finite. Thus, using
%\eqref{eq:one-frequency-Schatten} and the triangle inequality for
%an equivalent norm on \(S^{N,\infty}\) (recall that \(N>1\)),
%we obtain
%\begin{align*}
%\bigl\|
%[R_{{\rm Dunkl},j},M_f]
%\bigr\|_{S^{N,\infty}}
%&\lesssim
%\sum_{\sigma,\tau\in G}
%\sum_{u,v\in\mathbb Z^N}
%|c_{u,v}|
%\|\nabla f\|_{L^N}
%\\
%&\lesssim
%\|\nabla f\|_{L^N(\mathbb R^N)}.
%\end{align*}
%Indeed, the kernel of
%\([R_{{\rm Dunkl},j},M_f]\) is
%\[
%(f(y)-f(x))K_{{\rm Dunkl},j}(x,y),
%\]
%since our convention is
%\[
%[R_{{\rm Dunkl},j},M_f]
%=
%R_{{\rm Dunkl},j}M_f
%-
%M_fR_{{\rm Dunkl},j}.
%\]
%
%The preceding argument initially identifies the corresponding
%commutator form off the orbit diagonal \(d(x,y)=0\). Since that set
%has zero \(dw\times dw\)-measure, the NWO sum defines a unique
%operator in \(S^{N,\infty}\). By the standard local approximation
%of homogeneous Sobolev functions, or equivalently by first applying
%the argument to smooth Sobolev representatives and using the
%estimate above for differences, this operator is the closure of the
%usual commutator \([R_{{\rm Dunkl},j},M_f]\).
%
%Finally,
%\[
%\|\nabla f\|_{L^N(\mathbb R^N)}
%=
%\|f\|_{\dot W^{1,N}(\mathbb R^N)}.
%\]
%Hence
%\[
%[R_{{\rm Dunkl},j},M_f]
%\in
%S^{N,\infty}\bigl(L^2(\mathbb R^N,dw)\bigr)
%\]
%and
%\[
%\bigl\|
%[R_{{\rm Dunkl},j},M_f]
%\bigr\|_
%{S^{N,\infty}(L^2(\mathbb R^N,dw))}
%\lesssim
%\|f\|_{\dot W^{1,N}(\mathbb R^N)}.
%\]
%This proves \eqref{eq:Dunkl-critical-weak-Schatten-upper}.
%\end{proof}


\bibliographystyle{plain}
\bibliography{references}

\end{document}
